\documentclass[10pt,a4paper]{article}
\newcommand{\ReportShortTitle}{Polar decoding from the posterior}
\usepackage[T1]{fontenc}
\usepackage[utf8]{inputenc}
\usepackage{lmodern}
\usepackage{microtype}
\usepackage[a4paper,margin=24mm,headheight=14pt]{geometry}
\usepackage{amsmath,amssymb,mathtools}
\usepackage{booktabs,tabularx,array}
\usepackage{enumitem}
\usepackage{xcolor}
\usepackage{fancyhdr}
\usepackage{titlesec}
\usepackage[hidelinks]{hyperref}

\definecolor{ink}{HTML}{18212B}
\definecolor{accent}{HTML}{135F73}
\definecolor{muted}{HTML}{5E6A73}
\definecolor{pale}{HTML}{EEF5F7}
\color{ink}
\hypersetup{colorlinks=true,linkcolor=accent,urlcolor=accent}
\setlength{\parindent}{0pt}
\setlength{\parskip}{0.55em}
\setlist{nosep,leftmargin=1.55em}
\setcounter{secnumdepth}{2}
\titleformat{\section}{\Large\bfseries\color{accent}}{\thesection}{0.6em}{}
\titleformat{\subsection}{\large\bfseries\color{ink}}{\thesubsection}{0.55em}{}
\pagestyle{fancy}
\fancyhf{}
\fancyhead[L]{\small ODMA--URA research record}
\fancyhead[R]{\small\color{muted}\ReportShortTitle}
\fancyfoot[C]{\thepage}
\renewcommand{\headrulewidth}{0.3pt}
\newcolumntype{Y}{>{\raggedright\arraybackslash}X}
\newcommand{\R}{\mathbb{R}}
\newcommand{\C}{\mathbb{C}}
\newcommand{\Z}{\mathbb{Z}}
\newcommand{\F}{\mathbb{F}}
\newcommand{\one}{\mathbf{1}}
\newcommand{\given}{\,|\,}
\newcommand{\CN}{\mathcal{CN}}
\DeclareMathOperator{\diag}{diag}
\DeclareMathOperator{\Cov}{Cov}
\DeclareMathOperator{\Var}{Var}
\DeclareMathOperator{\supp}{supp}
\DeclareMathOperator*{\argminop}{arg\,min}
\DeclareMathOperator*{\argmaxop}{arg\,max}
\newcommand{\evidence}[1]{\textcolor{accent}{\textbf{Evidence status:}} #1}
\providecommand{\ReportUpdated}{30 August 2026}
\newcommand{\ProjectReport}[4]{%
  \pagenumbering{roman}
  \begin{titlepage}
    \vspace*{18mm}
    {\large\color{accent}\textbf{ODMA--URA research record}}\\[12mm]
    {\huge\bfseries #1}\\[5mm]
    {\Large\color{muted}#2}\\[18mm]
    \begin{tabularx}{\textwidth}{@{}p{0.16\textwidth}Y@{}}
      Stage & #3\\
      Updated & \ReportUpdated\\
      Status & #4
    \end{tabularx}
    \vfill
    \textcolor{muted}{This report is a chronological research record. Measured results, interpretations, and open
    hypotheses are labelled separately. Detailed numerical ledgers remain in \texttt{results/}.}
  \end{titlepage}
  \pagenumbering{arabic}
}

\hypersetup{citecolor=accent}
\newcommand{\E}{\mathbb E}
\newcommand{\GF}{\operatorname{GF}}
\newcommand{\calW}{\mathcal W}
\newcommand{\ind}{\mathbf 1}
\newcommand{\tr}{\operatorname{tr}}
\begin{document}
{\LARGE\bfseries Polar multiuser decoding from first principles}\par
{\large Exact multiset inference and a complete Bayes decoder}\par
{\color{muted}16 October 2026 --- Report 7; exact Bayesian foundations}

\textbf{Research question.} Given the fixed, deterministically balanced polar encoder of Report 8, what information
must an ideal receiver compute, and how can those calculations be completed exactly, even at prohibitive cost?
We begin with the physical model and the decision loss. D2 is a useful measured reference, not the
definition of either the problem or its solution.

\textbf{What this report establishes.} We derive the exact posterior in three equivalent representations, the
Bayes decision for multiplicity-aware PUPE, and a conditional polar-inference interface. We continue to exact
residual likelihoods, full count distributions, and a finite inside/outside decoder over polar message prefixes.
We also derive a decision certificate, finite Fourier and joint-separator evaluations, and an exact expansion
in connected groups of users. Further exact results identify sufficient continuation states, sharpen completion
and connected-remainder bounds, and locate the selected code's first same-pattern sum ambiguities.
Compatible elimination schedules, mixed-pattern state lower bounds and finite computational audits now test
how much these exact representations can actually compress. Overlap-parity contractions, reusable support
directions, adverse-event prior masses and a weighted-tree recurrence now provide concrete further reductions.
The three-user cycle is solved by streamed pair contractions, with full count probabilities and a precise
description of the dependence a later closure assumption would discard. A synthesis now connects exact merging,
reference choice, ambiguity and exact decision contrasts. This report contains exact identities, reductions,
obstructions and certified decisions. Algebraically motivated approximations and their error bounds have a
separate home in Report 9. Computational feasibility at research scale is still unresolved.

\section{Fix the experiment before deriving the receiver}
Use a synchronous, real, unit-gain Gaussian multiple-access channel:
\begin{equation}
 y=\sum_{k=1}^K\phi(w_k)+z,\qquad z\sim\mathcal N(0,\sigma^2 I_n),\qquad \|\phi(w)\|^2=1.
 \label{eq:model}
\end{equation}
The receiver knows $K\geq1$, $\sigma^2>0$ and the entire encoding rule. It knows neither the messages nor their
patterns. Initially $K<M$. There are no pilots, unknown fading coefficients or spatial observations in this model.
With the existing physical convention, $\sigma^2=(2B\,10^{E_b/N_0\,[{\rm dB}]/10})^{-1}$.

\begin{center}
\begin{tabularx}{\textwidth}{@{}lY@{}}\toprule
Symbol & Meaning and shape\\\midrule
$w=(p,v)$ & Pattern index $p\in\{0,\ldots,P-1\}$ and payload $v\in\GF(2)^{B_c}$\\
$P,V,M$ & $2^{B_p},2^{B_c},2^B=PV$; $B=B_p+B_c$\\
$L$ & $N_c\times B_c$ binary generator including CRC, frozen positions and polar transform\\
$b(v)$ & $(-1)^{Lv}\in\{-1,+1\}^{N_c}$, with the exponent evaluated modulo two\\
$S_p,D_p,A_p$ & Insertion $S_p\in\R^{n\times N_c}$; $D_p=\diag(\sqrt{e_p})$; $A_p=S_pD_p$\\
$\phi(p,v)$ & $A_pb(v)\in\R^n$; $S_p^TS_p=I$ and $\sum_i e_{p,i}=1$\\
$a$ & Unlabeled message counts in $\Z_{\geq0}^{M}$, with $\one^Ta=K$\\\bottomrule
\end{tabularx}
\end{center}

Freeze the code, placements $r(p,i)$ and the \emph{deterministic KL-balanced} energy bank from Report 8:
\[
 \min_{e\geq0}\ \sum_{p,i}\big[e_{p,i}\log(e_{p,i}/(1/N_c))-e_{p,i}+1/N_c\big],
 \quad \sum_i e_{p,i}=1,\quad \sum_{p,i:r(p,i)=r}e_{p,i}=P/n.
\]
Its converged matrix-scaling solution is the present reference. It uses no received data or trained parameters.
Feasibility and convergence must be checked. This construction balances resource energy; it does not prove optimal
PUPE, remove repeated polar-symbol functions, or settle the best placement. Those encoder questions remain parked
while we study decoding. The actual selected small-message code has $B_c=6,N_c=32$, CRC8 and $n=256$.

\clearpage
\section{Derive the common posterior}
IID uniform message draws are the default generative model at every $B$. Distinct sampling is the conditional
law given $D=\{w_i\ne w_j\text{ for all }i\ne j\}$, so
$\Pr_{\rm dist}(W\mid y)=\Pr_{\rm IID}(W\mid y,D)$.
Its prior difference has probability $1-(M)_K/M^K\leq K(K-1)/(2M)$, negligible when $K^2/M$ is small.
We retain both exact formulas so that conditioning is explicit; rare prior collisions need not be rare for every
individual observation.

\subsection{Counts: the shortest exact statement}
Let $\calW=\{0,\ldots,P-1\}\times\GF(2)^{B_c}$. Introduce the explicit matrix
\emph{as a mathematical definition}:
\begin{equation}
 \Phi=[\phi(w)]_{w\in\calW}\in\R^{n\times M},\qquad M=2^B,\qquad y=\Phi a+z.
\end{equation}
This notation does not require a practical receiver to store it. Gaussian noise gives the likelihood
$p(y\mid a)\propto\exp[-\|y-\Phi a\|^2/(2\sigma^2)]$. The prior depends on how transmissions are sampled:
\begin{align}
 \pi_{\rm dist}(a)&=\binom{M}{K}^{-1}\ind\{a\in\{0,1\}^M,\ \one^Ta=K\},\\
 \pi_{\rm iid}(a)&=\frac{K!}{M^K\prod_w a_w!}
                   \ind\{a\in\Z_{\geq0}^M,\ \one^Ta=K\}.
 \label{eq:priors}
\end{align}
For IID sampling, $K!/\prod_w a_w!$ is the number of ordered transmission sequences giving those counts.
For example, with two users, two different messages have twice the prior probability of two copies of one specified
message. A uniform prior over count vectors would describe a different experiment.

\begin{equation}
 \boxed{\Pr(a\mid y)=\frac{1}{Z_y}\pi(a)
          \exp\!\left[-\frac{\|y-\Phi a\|^2}{2\sigma^2}\right].}
 \label{eq:posterior}
\end{equation}
$Z_y$ is the sum of the numerator over all allowed $a$. This posterior is the shared information target.
Independent Bernoulli activities with mean $K/M$ do not impose exact cardinality; a fitted scalar offset imposing
$\sum_w\E[a_w]=K$ is not a substitute for \eqref{eq:priors}.

\subsection{Exchangeable slots: expose the code structure}
Write $W=(w_1,\ldots,w_K)$, where slots are mathematical labels, not user identities. The slot prior is $M^{-K}$ for
IID sampling and $(M)_K^{-1}\ind\{w_i\ne w_j\text{ for }i\ne j\}$ for distinct sampling, where
$(M)_K=M(M-1)\cdots(M-K+1)$. Then
\begin{equation}
 \Pr(W\mid y)\propto\pi(W)
       \exp\!\left[-\frac{\|y-\sum_k A_{p_k}b(v_k)\|^2}{2\sigma^2}\right].
 \label{eq:slots}
\end{equation}
Summing \eqref{eq:slots} over all orderings with counts $a$ gives \eqref{eq:posterior}, including the factorial
weight. The representations are equivalent; their \emph{modes} need not be. Maximizing a single IID slot assignment
is maximum likelihood, whereas IID count MAP includes its multiplicity weight.

To see the actual coupling, put $x_k=A_{p_k}b(v_k)$ and expand:
\begin{align}
 \left\|y-\sum_kx_k\right\|^2
 &=\|y\|^2+K-2\sum_k y^Tx_k+2\sum_{k<j}x_k^Tx_j,\\
 \Pr(W\mid y)&\propto\pi(W)
   \exp\!\left[\frac{\sum_k y^Tx_k-\sum_{k<j}x_k^Tx_j}{\sigma^2}\right].
 \label{eq:interactions}
\end{align}
The first sum is individual evidence; the second accounts for messages explaining the same observations.
Taking the largest isolated correlations discards the second sum. Balancing an average frame operator does not
make these pairwise interactions vanish.

\textbf{Symmetry matters.} The exact posterior is invariant to permuting slots. In particular,
$\E[a_w\mid y]=K\Pr(w_1=w\mid y)$. Slot marginals do not each identify a different physical user. Separation into
different slot beliefs is an approximation or a conditional hypothesis, not an exact property of the posterior.

\clearpage
\section{What does an optimal decoder optimize?}
\subsection{Multiplicity-aware PUPE and its exact optimal decision}
The primary objective is recovery of transmitted message \emph{instances}. The receiver returns a multiset
$\widehat a\in\mathbb Z_+^M$ with at most $K$ occurrences. One returned copy can match only one transmitted copy:
\begin{equation}
 \boxed{\ell_{\rm multi}(a,\widehat a)
 =1-\frac1K\sum_w\min(a_w,\widehat a_w).}
 \label{eq:loss}
\end{equation}
When both totals equal $K$, this is $\|a-\widehat a\|_1/(2K)$. For a shorter output use the first expression:
missing copies still count as errors. For truth $(u,u,v)$, output $(u,v)$ has error $1/3$ and $(u,u,v)$ has zero
error. Distinct-message truth reduces to ordinary message-list PUPE.

Define posterior count tails $T_{w,j}(y)=\Pr(a_w\geq j\mid y)$ for $j=1,\ldots,K$. The elementary identity
$\min(a_w,d)=\sum_{j=1}^d\ind\{a_w\geq j\}$ gives
\begin{equation}
 \E[\ell_{\rm multi}(a,\widehat a)\mid y]
 =1-\frac1K\sum_w\sum_{j=1}^{\widehat a_w}T_{w,j}(y).
 \label{eq:bayes}
\end{equation}
An exact Bayes action allocates $K$ copies by repeatedly choosing the largest next available tail
$T_{w,\widehat a_w+1}$ and incrementing that message's count. This greedy allocation is globally optimal:
each message has nonincreasing tails, so choosing its $j$th copy never requires skipping a more valuable
earlier copy. Ties can be resolved while preserving these prefixes. Filling all $K$ positions cannot increase
this loss. This specifies an ideal decision, not an efficient search across the $M$ messages.

Under distinct sampling only $T_{w,1}=\Pr(a_w=1\mid y)$ is nonzero. The rule becomes top-$K$ inclusion probability.
Under IID sampling posterior means $\eta_w=\E[a_w\mid y]=\sum_jT_{w,j}$ are generally insufficient: a large
expected count does not specify the separate values of returning one or two copies.

\subsection{MAP is a different decision}
Count MAP minimizes $\ind\{\widehat a\ne a\}$, the error on the whole multiset. It need not minimize
\eqref{eq:loss}. For example, consider this normalized posterior on pairs of four distinct messages:
\[
 \Pr(01,02,03,12,13,23\mid y)=(.30,.22,.01,.24,.01,.22).
\]
MAP selects $01$. The inclusion probabilities are $(.53,.55,.68,.24)$, so the Bayes action is $12$.
Their conditional PUPE risks are $0.460$ and $0.385$, respectively. This is a decision-theoretic example, not a new
measured operating point. An optimization route aimed at MAP must acknowledge this change of decision target.

\subsection{Metric correction and historical comparisons}
Before the 11 October 2026 correction, the main array/count and matched benchmark scorers used
\[
 \ell_{\rm list}(a,\mathcal L)=1-\frac1K\sum_{w\in\mathcal L}a_w.
\]
This credits every copy when a message appears once. Its Bayes list ranks $\eta_w$, but it was not the intended
objective. It remains as \texttt{list\_pupe}; new primary scores have \texttt{pupe\_metric\_version=2}.
Old files are unchanged. Distinct-message results with unchanged valid output support retain their PUPE;
historical IID results require saved hard predictions or checkpoint re-evaluation to establish multiplicity recovery.

The matched adapter now preserves hard-count copies. Count-array readout caps occurrences at $K$ by descending
estimated count with stable index ties; already valid outputs are unchanged. This declared budget policy is not
the Bayes allocation above. Raw count errors, overflow and false alarms remain diagnostics.
The benchmark's \texttt{strict\_pupe} still treats every duplicated-message user as erroneous.

The decoder investigation retains posterior count tails and the Bayes multiset decision in \eqref{eq:bayes}
as its target. Approximating the calculation does not change this decision objective.

\clearpage
\section{The first exact polar operation}
\subsection{Condition on the other messages, without pretending they are known}
For a particular slot $k$, fix a \emph{hypothesized} assignment $W_{-k}$ and subtract its waveform:
\begin{equation}
 r_{-k}=y-\sum_{j\ne k}\phi(w_j).
\end{equation}
This is an operation inside inference; supplying the true $W_{-k}$ would be an oracle diagnostic. Let
$\mathcal E_k$ contain all messages in the IID case and exclude $W_{-k}$ in the distinct case. Unit energy gives
\begin{align}
 \Pr(w_k=(p,v)\mid y,W_{-k})
 &\propto\ind\{(p,v)\in\mathcal E_k\}
          \exp[-\|r_{-k}-A_pb(v)\|^2/(2\sigma^2)]\\
 &\propto\ind\{(p,v)\in\mathcal E_k\}
          \exp[\tfrac12\lambda_p^Tb(v)],\qquad
 \boxed{\lambda_p=\frac{2A_p^Tr_{-k}}{\sigma^2}.}
 \label{eq:conditional}
\end{align}
The squared candidate energy is one, so it disappears from comparisons across both payloads and patterns.
At an occupied coordinate the LLR is $2\sqrt{e_{p,i}}\,r_{-k,r(p,i)}/\sigma^2$. Polar constraints enter through
$b(v)$, not by changing the physical variance.
For $K=1$, $r_{-k}=y$: this is already the full posterior, with no multiuser approximation.

\subsection{Three different operations, often mistakenly conflated}
For any supplied $\lambda$, define
\begin{equation}
 Z_{\rm pol}(\lambda)=\sum_{v\in\GF(2)^{B_c}}\exp[\lambda^Tb(v)/2].
 \label{eq:zpol}
\end{equation}
\begin{enumerate}
\item \textbf{Maximize:} find the valid payload with largest $\lambda^Tb(v)/2$.
      This is the conditional mode or a candidate-search operation.
\item \textbf{Sum:} evaluate $Z_{\rm pol}$ and marginals; for instance
      $\E[b_i\mid\lambda]=2\,\partial\log Z_{\rm pol}/\partial\lambda_i$.
      Pattern probabilities require these sums.
\item \textbf{Sample:} draw a payload from the distribution proportional to that exponential.
      This supports posterior sampling, not just optimization.
\end{enumerate}
With uniform messages and no exclusion, $\Pr(p_k=p\mid y,W_{-k})=Z_{\rm pol}(\lambda_p)/\sum_qZ_{\rm pol}(\lambda_q)$.
For distinct sampling, remove forbidden full-message terms from each pattern's sum.
\emph{Choosing the pattern with the strongest single payload is not marginal pattern inference.}
If a decoder uses normalized bit-channel probabilities, their pattern-dependent normalization constants must be
restored before comparing different patterns. Using the unnormalized scores in \eqref{eq:zpol} avoids that ambiguity.

Polar encoding has a compact XOR circuit. SC and bounded-list SCL exploit that circuit efficiently
\cite{arikan,tal}, but a finite list does not guarantee the exact maximum, partition sum or sample in
\eqref{eq:zpol}. In particular, future frozen/CRC constraints and pruning cannot be ignored when claiming exact
posterior sampling. The exact conditional interface is established; a cheap exact implementation is not.

\subsection{Continue the exact sum before choosing an approximation}
The unconditional marginal is
$\Pr(w_k\mid y)=\sum_{W_{-k}}\Pr(w_k\mid y,W_{-k})\Pr(W_{-k}\mid y)$.
Replacing uncertain interference by one subtraction is a substantive approximation to this sum.
Exact enumeration requires $\binom{M}{K}$ distinct or $\binom{M+K-1}{K}$ IID count states; compact slot notation still
allows $M^K$ assignments. Polar structure compresses the description, not these sums automatically.

Nothing in this identity requires choosing an approximation now. We next turn the sum into complete finite
procedures and ask which states can be combined without losing the count tails needed by the decision.
Approximation assumptions are developed separately in Report 9 from the bottlenecks of these exact procedures.

\clearpage
\section{Complete the exact residual-likelihood decoder}
\label{sec:residual-exact}
\subsection{One function for every possible explanation of a residual}
For an allowed full-message alphabet $\mathcal S\subseteq\calW$, define
\begin{align}
 g(r)&=\exp[-\|r\|^2/(2\sigma^2)],\\
 H_m^{\mathcal S}(r)&=
 \sum_{(w_1,\ldots,w_m)\in\mathcal S^m}
       g\!\left(r-\sum_{i=1}^m\phi(w_i)\right),\qquad H_0^{\mathcal S}(r)=g(r).
 \label{eq:ordered-evidence}
\end{align}
$H_m$ is a sum of likelihoods, not a normalized density. Multiplication by
$|\mathcal S|^{-m}(2\pi\sigma^2)^{-n/2}$ gives the observation density for $m$ IID uniform draws from $\mathcal S$.
Keep the unnormalized definition when comparing different alphabets or different $m$.

Separate the first message from its remaining $m-1$ companions:
\begin{equation}
 \boxed{H_m^{\mathcal S}(r)=
        \sum_{w\in\mathcal S}H_{m-1}^{\mathcal S}(r-\phi(w)).}
 \label{eq:residual-recursion}
\end{equation}
Each branch subtracts a \emph{hypothesis}; summing all branches integrates out the interference exactly.
The base case contains the physical noise once. Adding noise variance or a new Gaussian likelihood at every
recursion level would change the model.

For a small example, $H_2(r)=\sum_u\sum_v g(r-\phi(u)-\phi(v))$. Different orderings of $u,v$ count separately,
as the IID slot prior requires. Repeated messages are allowed, as are different payloads on the same pattern.
Every pattern is included in the alphabet; no true pattern is supplied.

\subsection{Connect the base operation to the polar partition function}
With the full alphabet and unit energy,
\begin{equation}
 H_1^{\calW}(r)=
 e^{-(\|r\|^2+1)/(2\sigma^2)}
 \sum_{p=0}^{P-1}Z_{\rm pol}\!\left(\frac{2A_p^Tr}{\sigma^2}\right).
 \label{eq:h1-polar}
\end{equation}
This follows by expanding $\|r-A_pb(v)\|^2$ and then summing valid payloads. Thus the exact polar evidence
operation is the one-user boundary of an exact multiuser recursion.
For an excluded alphabet, sum only its allowed payloads. Positive enumeration avoids cancellation from subtracting
nearly equal partition sums.

The recursion is already a terminating algorithm: recursively generate every valid message, evaluate the Gaussian
at depth zero, and add. Naively there are $M^m$ leaves. Reusing identical subproblems can help, but writing a
recurrence does not by itself reduce this exponential count. In particular, $H_m$ is a function on $\R^n$;
it is not a table of only $m+1$ numbers.

\textbf{What has improved?} Unknown interference is now a specified, evaluable function rather than an
unspecified marginalization. The next step is to obtain the \emph{right} queries of that function for multiset PUPE.

\clearpage
\subsection{From residual evidence to every posterior count probability}
Fix a full message $w$ and write $q_{w,j}=\Pr(a_w=j\mid y)$. Exactly $j$ of the $K$ labeled slots can contain $w$
in $\binom Kj$ ways. The remaining slots must exclude $w$. Therefore
\begin{equation}
 \boxed{q_{w,j}=
 \binom Kj\,
 \frac{H_{K-j}^{\calW\setminus\{w\}}(y-j\phi(w))}
      {H_K^{\calW}(y)},\qquad j=0,\ldots,K.}
 \label{eq:count-probability}
\end{equation}
All terms are nonnegative. This is not an independence approximation: the residual function still sums the
joint explanations of all remaining transmissions. As checks,
$\sum_jq_{w,j}=1$ and $\sum_w\sum_j j q_{w,j}=K$.

The exclusion is essential. If we instead allow the remaining slots to contain $w$, we obtain
\begin{equation}
 \E\!\left[\binom{a_w}{j}\middle|y\right]=
 \binom Kj\,\frac{H_{K-j}^{\calW}(y-j\phi(w))}{H_K^{\calW}(y)}.
 \label{eq:factorial-moment}
\end{equation}
This counts choices of $j$ marked copies, including states with more than $j$ copies.
At $j=1$ it gives $\E[a_w\mid y]$, the quantity in the shared discussion.
That mean alone is insufficient for the corrected loss. Instead form
\[
 T_{w,j}=\sum_{t=j}^Kq_{w,t}
\]
and use the allocation in \eqref{eq:bayes}. The physical model, exact inference, and optimal decision are now connected.

\subsection{Distinct messages require a different recursion}
Let $D_m^{\mathcal S}$ sum the same Gaussian weight over ordered tuples with distinct entries. Then
\begin{equation}
 D_0^{\mathcal S}=g,\qquad
 D_m^{\mathcal S}(r)=\sum_{w\in\mathcal S}
 D_{m-1}^{\mathcal S\setminus\{w\}}(r-\phi(w)).
 \label{eq:distinct-recursion}
\end{equation}
Set $D_m^{\mathcal S}=0$ when $m>|\mathcal S|$. Here
$q_{w,1}=K D_{K-1}^{\calW\setminus\{w\}}(y-\phi(w))/D_K^{\calW}(y)$,
$q_{w,0}=1-q_{w,1}$ and higher counts have zero probability.
This is the conditioned source model, not an IID recursion with a different final threshold.

\subsection{A complete finite decoder, with its guarantee}
\begin{enumerate}
\item Generate valid $(p,v)$ messages using the fixed polar/CRC encoding rule.
\item Compute $H_K(y)$ and the excluded-alphabet functions in \eqref{eq:count-probability}
      (or their distinct counterparts).
\item Normalize to $q_{w,j}$ for every $w,j$ and sum to obtain count tails.
\item Allocate $K$ occurrences by the largest next tail, preserving each message's prefix of copies.
\end{enumerate}
All sums are finite, and $H_K(y)>0$ for $\sigma^2>0$. In exact arithmetic this algorithm returns a global
minimum of posterior expected multiset PUPE for every observation. Integrating that conditional guarantee over
$y$ gives minimum average PUPE under the stipulated prior. It promises neither zero error nor fast computation.

For completeness, the same evidence functions also give exact sequential posterior sampling:
choose a first slot with probability $H_{K-1}(y-\phi(w))/H_K(y)$, subtract it, and repeat with one fewer slot.
These probabilities telescope to the full slot posterior. This finite construction needs exact evidence;
it is stronger than a stationary-target statement for a finite Monte Carlo run, but just as potentially expensive.

\clearpage
\section{An exact inside/outside decoder over message prefixes}
\label{sec:tree-exact}
\subsection{Replace ordered tuples by positive signal-sum measures}
The preceding algorithm repeatedly counts permutations. We can remove that repetition and share computations
between message queries without changing the posterior.
For a message subset $\mathcal S$, define the discrete measure
\begin{equation}
 \mu_{\mathcal S,k}
 =\sum_{\substack{a:\,\supp(a)\subseteq\mathcal S\\\one^Ta=k}}
       \beta(a)\,\delta_{\Phi a},
 \qquad
 \beta(a)=\prod_w\frac1{a_w!}\quad\text{(IID)}.
 \label{eq:measure}
\end{equation}
$\delta_s$ means a point mass at waveform sum $s$. For distinct sampling restrict $a_w$ to $0,1$, so
$\beta(a)=1$. These are unnormalized measures. Their total masses are
$|\mathcal S|^k/k!$ for IID and $\binom{|\mathcal S|}{k}$ for distinct sampling.
For the empty alphabet only $\mu_{\varnothing,0}=\delta_0$ is nonzero.

For two measures, convolution adds their signals and multiplies their weights:
$\delta_s*\delta_t=\delta_{s+t}$. If $\mathcal S=\mathcal A\mathbin{\dot\cup}\mathcal B$ is a disjoint union,
\begin{equation}
 \boxed{\mu_{\mathcal S,k}=
        \sum_{q=0}^k\mu_{\mathcal A,q}*\mu_{\mathcal B,k-q}.}
 \label{eq:measure-convolution}
\end{equation}
Proof: every count vector has a unique restriction to the two disjoint alphabets.
Its total count and waveform add, and $\beta(a)$ factors. There is no binomial coefficient here:
the leaf factorials already implement the IID prior.

Now partition complete messages by their binary prefixes. The first $B_p$ levels select the pattern;
the next $B_c$ levels select the payload. Each leaf is exactly one valid full message, whose waveform is
constructed by polar encoding. The split is a partition of \emph{possible messages}, not of physical users.
Unknown pattern occupancies and arbitrary same-pattern multiplicities remain in the sum.

\subsection{The inside pass}
For a node $u$, let $\mathcal S_u$ be its leaf alphabet and $I_u(k)=\mu_{\mathcal S_u,k}$.
At the leaf for message $w$:
\begin{equation}
 I_w(k)=\frac1{k!}\delta_{k\phi(w)}\quad\text{(IID)},\qquad
 I_w(0)=\delta_0,\ I_w(1)=\delta_{\phi(w)},\ I_w(k\geq2)=0\quad\text{(distinct)}.
 \label{eq:leaf}
\end{equation}
At each parent apply \eqref{eq:measure-convolution} to its two children, keeping counts $0,\ldots,K$.
The root yields the evidence
\begin{equation}
 \mathcal Z(y)=\int g(y-s)\,I_{\rm root}(K)(ds)
 =H_K^{\calW}(y)/K!\quad\text{(IID)}
 \label{eq:tree-evidence}
\end{equation}
and $D_K^{\calW}(y)/K!$ for distinct sampling. The omitted prior normalization is constant over all states
of total $K$ and cancels from posterior ratios. Apply the Gaussian after summing signals; multiplying
Gaussian scores of the child signals would incorrectly treat their shared observation as independent.

\clearpage
\subsection{The outside pass answers the missing count queries}
Define $O_u(k)=\mu_{\calW\setminus\mathcal S_u,k}$, the measure for messages outside node $u$.
Initialize $O_{\rm root}(0)=\delta_0$ and $O_{\rm root}(k>0)=0$.
If $u$ has parent $t$ and sibling $v$, its complement is the disjoint union of the parent's complement and $v$:
\begin{equation}
 \boxed{O_u(k)=\sum_{q=0}^k O_t(q)*I_v(k-q).}
 \label{eq:outside}
\end{equation}
A pass from root to leaves therefore computes the complementary measures for every message.
At a leaf,
\begin{equation}
 \boxed{q_{w,j}=\frac{\beta_w(j)}{\mathcal Z(y)}
       \int g(y-j\phi(w)-t)\,O_w(K-j)(dt),}
 \label{eq:tree-query}
\end{equation}
where $\beta_w(j)=1/j!$ in the IID case and is $1$ for $j=0,1$, $0$ otherwise in the distinct case.
For IID sampling, substituting $O_w(m)=\mu_{\calW\setminus\{w\},m}$ and
$H_m=m!\int g\,d\mu$ recovers \eqref{eq:count-probability}, including its binomial coefficient.

\textbf{The full algorithm is now explicit:} build the prefix tree, compute all inside measures, propagate
outside measures, evaluate \eqref{eq:tree-query}, sum count tails, and allocate the $K$ decoded occurrences.
No candidate proposer, noise approximation, true pattern, training, or unspecified marginal oracle is required.
The tiny implementation in \path{tests/polar_exact_inference_test.py} executes these steps.

\subsection{What may be compressed, and what must be retained}
If two configurations produce exactly the same sum $s$, their weights may be added in an inside or outside
measure: every subsequent convolution and Gaussian evaluation uses only that sum.
This is an exact local compression. However, the root sum distribution alone does not determine $q_{w,j}$.
Different configurations with the same signal can have different message counts. The outside pass keeps the
query's own leaf separate, so its count information survives the aggregation.

For example, if messages $u$ and $v$ have the same waveform, observing that waveform does not reveal which was
sent. Summing their evidence is valid; assigning the combined evidence to one arbitrarily chosen message is not.
The certificates deliberately include this ambiguity to check the distinction.

\subsection{Where the computational burden has moved}
Let $s_u(k)$ be the number of distinct stored sum states at node $u$. Combining two children costs, before
coalescing equal states,
$\sum_{k=0}^K\sum_{q=0}^k s_{\rm left}(q)s_{\rm right}(k-q)$ pair additions,
each also requiring signal-key work. Outside passes have the same kind of cost.
Without useful coincidences, $s_u(k)$ can be $\binom{|\mathcal S_u|+k-1}{k}$ for IID sampling.
Thus this construction removes slot permutations and reuses shared sums; it does not establish polynomial
complexity or sub-exhaustive payload search.

The full tree has $M$ leaves. It can generate waveforms procedurally and avoid storing the dense $n\times M$
matrix, while still having exponential dependence on $B$. This is a complete exact decoder and a reference
for future reductions, not a claim of scalable execution. Small certificates use floating point arithmetic;
large evidence calculations require scaled/log-domain sums, and rigorous numerical bounds require controlled rounding.

\clearpage
\section{Use polar algebra without losing association}
\label{sec:polar-exact}
\subsection{Group identical parity functions exactly}
Write the $i$th row of $L$ as $\ell_i\in\GF(2)^{B_c}$, and define the binary characters
$\chi_g(v)=(-1)^{g^Tv}$. Then $b_i(v)=\chi_{\ell_i}(v)$.
For the distinct row set $\mathcal G$, collect
\[
 C_{p,g}=\sum_{i:\ell_i=g}\sqrt{e_{p,i}}\,S_p[:,i],
 \qquad
 \phi(p,v)=\sum_{g\in\mathcal G}C_{p,g}\chi_g(v).
\]
For counts $a_{p,v}$ the total waveform is therefore
\begin{equation}
 s(a)=\sum_{p,g}C_{p,g}\,t_{p,g},\qquad
 t_{p,g}=\sum_v a_{p,v}\chi_g(v)\in\{-K,\ldots,K\}.
 \label{eq:integer-signature}
\end{equation}
Use the integer array $t$ as the convolution key and the known $C$ only when evaluating the physical signal.
Addition of keys is exact; no tolerance-based merging of real waveforms is needed. Equal signatures can be
combined even when their underlying count vectors differ, with the outside-query safeguard above.
Different signatures may also produce the same physical signal; retaining them separately is redundant but correct.

The selected $32\times6$ generator has 29 distinct rows, including the zero character. That is a small exact
reduction of the signature dimension, not a disappearance of its state explosion.
If all $2^{B_c}$ characters were retained, the Walsh inverse would recover each pattern's counts:
\[
 a_{p,v}=2^{-B_c}\sum_{g\in\GF(2)^{B_c}}t_{p,g}\chi_g(v).
\]
Only a subset is transmitted here. General count recovery cannot assume that this incomplete transform is
invertible; restricted-$K$ identifiability is a separate question.

\subsection{An exact one-user transform, with a precise limit}
For supplied LLRs, group $h_g=\tfrac12\sum_{i:\ell_i=g}\lambda_i$ and put $h_g=0$ for missing characters.
Then every payload log score is
\[
 f(v)=\tfrac12\lambda^Tb(v)=\sum_g h_g\chi_g(v),\qquad
 Z_{\rm pol}(\lambda)=\sum_v e^{f(v)}.
\]
The unnormalized Walsh transform computes all $f(v)$ in $O(B_c2^{B_c}+N_c)$ operations rather than
$O(N_c2^{B_c})$ separate codeword scores. This is an exact finite-alphabet improvement; it still enumerates
the payload alphabet. Exponentiation and the multiuser interactions are not diagonalized by this transform.

\subsection{Why an aggregate XOR butterfly needs more state}
For two users, incoming bit pairs $(0,0),(1,1)$ and $(0,1),(1,0)$ both have one one in each input column.
Their per-user XOR outputs are respectively $(0,0)$ and $(1,1)$.
Thus counts of ones on the two input wires do not determine the output count: we also need their
joint pairing across users. A polar butterfly on marginal counts alone is not an exact recursion.
Full joint separator states, or an equivalent lossless association representation, could be propagated instead,
but their size must be derived. The prefix convolution above avoids this unsupported simplification by
constructing valid complete-codeword signatures before combining them.

\clearpage
\subsection{An equivalent exact polar and pattern factorization}\label{sec:bit-factorization}
For slot $k$, introduce payload bits $v_k$, polar input bits $u_k$, and intermediate bits at all butterfly stages.
The local polar kernel has the indicator
\begin{equation}
 f_{\rm ker}(a,b,c,d)=\ind\{c=a\oplus b,\ d=b\}.
\end{equation}
Frozen input bits equal zero. Information positions contain the payload and its deterministic CRC; the CRC can be
represented by its binary linear constraints. Let $\mathcal P_k$ be the product of these indicators, including
payload insertion. It allows exactly one internal encoding path for each $v_k$.

Introduce real-valued variables $x_{k,r}$ supported on their finite encoder-induced alphabets, with mapping factors
\begin{equation}
 \mathcal M_k=\prod_{r=1}^{n}
 \ind\!\left\{x_{k,r}=\sum_{i=1}^{N_c}
       \ind\{r(p_k,i)=r\}\sqrt{e_{p_k,i}}\,(-1)^{c_{k,i}}\right\}.
\end{equation}
There is \emph{one shared pattern variable} $p_k$ across all resources of a slot. Replacing it by independent on/off
choices at each resource would admit invalid patterns. Conditional on $p_k$, each resource has at most one local
symbol from that slot; different slots may use the same pattern.

\begin{equation}
 \Pr(p,v,u,c,x,\ldots\mid y)\propto
 \pi(W)\prod_k\mathcal P_k\mathcal M_k
 \prod_{r=1}^{n}\exp\!\left[-\frac{(y_r-\sum_kx_{k,r})^2}{2\sigma^2}\right].
 \label{eq:graph}
\end{equation}
For distinct sampling, retain the all-different constraint on \emph{full} messages in $\pi(W)$.
For IID sampling it is absent. Summing out the deterministic internal variables gives \eqref{eq:slots} exactly.
Thus counts, slots and the bit/pattern graph describe the same distribution before approximation.

Every resource factor has a finite joint sum, and exact variable elimination recovers the same posterior.
The local XOR circuit does not bound the size of intermediate joint factors. In particular, ordinary scalar
messages on the loopy graph are not an additional exact decoder. Section~\ref{sec:elimination-exact} treats
exact elimination; assumptions about messages or interference belong to the separate approximation report.

\clearpage
\section{What an exact decoder may safely forget}
\label{sec:exact-reductions}
\subsection{First remove observation directions that cannot distinguish messages}
Let $\mu=M^{-1}\sum_w\phi(w)$ and
$V=\operatorname{span}\{\phi(w)-\mu:w\in\calW\}$. Choose an orthonormal basis matrix $Q$ for $V$, and write
$P_V=QQ^T$. Every allowed sum of $K$ messages lies in $K\mu+V$, so
\begin{equation}
 \|y-\Phi a\|^2
 =\|Q^T(y-K\mu)-Q^T(\Phi a-K\mu)\|^2
  +\|(I-P_V)(y-K\mu)\|^2.
 \label{eq:sufficient-observation}
\end{equation}
The last term is identical for every hypothesis. It cancels from every posterior ratio. Thus
$Q^T(y-K\mu)$ is an exactly sufficient observation for this decoder. Its noise covariance is still
$\sigma^2 I$; no interference variance or Gaussian approximation has been introduced.

The space can be found without enumerating the full payload alphabet. Character orthogonality gives
\begin{equation}
 \mu=\frac1P\sum_p C_{p,0},\qquad
 \Sigma=\frac1P\sum_{p,g}C_{p,g}C_{p,g}^T-\mu\mu^T,\qquad V=\operatorname{range}\Sigma.
 \label{eq:variation-moments}
\end{equation}
Use $C_{p,0}=0$ if the zero character is absent. To prove the last equality, note that
$d^T\Sigma d=\operatorname{Var}_w(d^T\phi(w))$ vanishes exactly when $d^T\phi(w)$ is constant over all messages.
This is the same encoder moment used in Report 8, now interpreted as a sufficient-statistic calculation.

\subsection{The actual bank rules out a global dimension reduction}
For the selected code on one \emph{known} pattern, the 28 nonzero character vectors have disjoint nonempty
supports, so they are linearly independent. At known occupancy, the signal variation therefore has dimension
28, rather than 32. The zero character contributes a known offset; repeated characters contribute no new
directions. If occupancy is unknown, that offset must instead be retained as a state variable.

For the full $B=12/14$, $n=256$ banks and seeds 3201--3203, however, $V=\R^{256}$.
This has a structural certificate, not just a numerical rank estimate:
\begin{enumerate}
\item Each nonzero $C_{p,g}$ belongs to $V$, since it is the character-weighted average of the pattern's
      codewords, with coefficients summing to zero.
\item A character occurring at one symbol isolates a coordinate axis. In a repeated-character vector, subtract
      already isolated axes; if one support coordinate remains, that axis also belongs to $V$.
\item This isolates 254, 254, 253 axes at $B=12$ and 255 at each $B=14$ seed. The only constant symbol is
      position zero in the local code. Differences between pattern means, anchored at an already isolated
      mean coordinate, isolate the remaining axes in all six banks.
\end{enumerate}
All used coefficients are positive, so no numerical threshold is required. Consequently the global projection
above saves no dimensions here. It remains useful after conditioning or at smaller alphabet cuts.

\clearpage
\subsection{A stronger rule at a fixed cut: only preserve interaction with the future}
Split the message alphabet into disjoint sets $A,B$. Fix $q$ messages in $A$ and $h=K-q$ in $B$.
Let $\mathcal T_{B,h}$ be the allowed future sums, and choose any point $t_0$ in their affine hull. Define
\[
 V_{B,h}=\operatorname{span}\{t-t_0:t\in\mathcal T_{B,h}\},\qquad P_{B,h}=Q_{B,h}Q_{B,h}^T.
\]
For a partial sum $s$ from $A$ and a completion $t$ from $B$, orthogonal decomposition gives
\begin{align}
 g(y-s-t)
 &=\eta_h(s)\,
 \exp\!\left[-\frac{\|Q_{B,h}^T(y-t_0-s)-Q_{B,h}^T(t-t_0)\|^2}{2\sigma^2}\right],
 \label{eq:continuation-projection}\\
 \eta_h(s)&=\exp\!\left[-\frac{\|(I-P_{B,h})(y-t_0-s)\|^2}{2\sigma^2}\right].
 \nonumber
\end{align}
Thus all partial states with the same $Q_{B,h}^Ts$ can be merged after multiplying their weights by
$\eta_h(s)$. Every later completion sees precisely the same remaining likelihood. For a count query, merge
the coefficients of its marker polynomial, not just the scalar weights. This retains the count information
required by the loss even when message identities have different past explanations.

This criterion is also necessary for merging two signal states \emph{up to a scalar weight} while preserving
each allowed continuation. Indeed
\[
 \log\frac{g(y-s-t)}{g(y-s'-t)}
 =\frac{(s-s')^T(y-t)-(\|s\|^2-\|s'\|^2)/2}{\sigma^2}
\]
is independent of $t$ if and only if $s-s'\perp V_{B,h}$.
The conclusion concerns proportional likelihood rows with the same future feasibility; it is not a universal
minimality theorem for every algebraic representation.

No multiuser enumeration is needed merely to find this space. For IID sampling and $h\geq1$ it equals
$\operatorname{span}\{\phi(u)-\phi(v):u,v\in B\}$: compare completions differing in one user.
For distinct sampling the same proof works when $1\leq h<|B|$, by fixing $h-1$ other messages outside
$\{u,v\}$. For $h=0$, or distinct $h=|B|$, there is just one future sum and the space is zero.
For nonempty $B$, use $t_0=h|B|^{-1}\sum_{w\in B}\phi(w)$, the mean of the admissible future sums.
If $B$ is empty, only $h=0$ is feasible and $t_0=0$.

\textbf{Boundary of the rule.} Disjoint alphabet cuts make the prior and distinct exclusions compatible with this
merging. Two labeled histories with different used-message sets generally have different legal futures; projection
equality alone cannot merge them. The weights also depend on the cut, occupancy and observation. Equation
\eqref{eq:continuation-projection} is an exact cut calculation, not permission to multiply the Gaussian likelihood
again at every recursion level. A complete elimination order must consume each likelihood factor once.

\clearpage
\subsection{A small geometric separator need not be a small table or tensor}
If no future message touches a resource, that resource lies outside $V_{B,h}$ and its likelihood may be absorbed
into the current weight. This gives a resource-frontier interpretation of the projection rule. More general
linear relations can remove additional directions even when supports overlap.

The actual polar payload prefixes provide a concrete local profile. Once a pattern and $\ell$ leading payload
bits are fixed, restrict each character row to its remaining $6-\ell$ bits. Equal restricted rows merge;
zero rows contribute only a known offset. The merged vectors still have disjoint supports, so the number of
distinct nonzero restricted rows is the exact one-message variation dimension. For $\ell=0,\ldots,6$ it is
\[
 28,\quad23,\quad15,\quad7,\quad3,\quad1,\quad0.
\]
For IID positive occupancy these are also the future-sum variation dimensions. Distinct sampling has the
same dimensions at interior occupancies, and dimension zero when all remaining messages must be selected.
This establishes useful local reductions, while leaving the number of distinct sum states to be determined.

At $\ell\geq2$, every character on the remaining bits is present. Write $V'=2^{6-\ell}$.
The full Walsh inverse therefore recovers the restricted payload counts from the local sum and known occupancy
$h$. There are exactly $\binom{V'+h-1}{h}$ different local sum states under IID sampling and $\binom{V'}h$
under distinct sampling. For example, $V'=16,h=4$ gives 3,876 IID or 1,820 distinct states.
Equal-signal merging cannot reduce these local banks. This statement concerns the full local signal, or its
projection onto its own variation space; a projection against a different future alphabet is a separate question.

There are nevertheless two different size questions:
\begin{enumerate}
\item How many coordinates describe the retained projection?
\item How many distinct values, or linearly independent likelihood rows, occur in those coordinates?
\end{enumerate}
The first can be small while the second remains large. A one-dimensional example proves the point.
Suppose partial and future sums take scalar values $i,j\in\{0,\ldots,d-1\}$. Expanding the Gaussian gives
\[
 e^{-(y-i-j)^2/(2\sigma^2)}
 = e^{-y^2/(2\sigma^2)}
   e^{(yi-i^2/2)/\sigma^2}e^{(yj-j^2/2)/\sigma^2}
   \underbrace{e^{-ij/\sigma^2}}_{E_{ij}}.
\]
The first three factors are a constant and nonzero row/column scalings. They do not change rank.
Put $x_j=e^{-j/\sigma^2}$. Then $E_{ij}=x_j^i$ is a Vandermonde matrix, with
\[
 \det E=\prod_{0\leq j<\ell<d}(x_\ell-x_j)\ne0
 \qquad(\sigma^2>0).
\]
Its rank is $d$, however large $d$ becomes, even though the shared signal coordinate is one-dimensional.
Consequently a rank-one bilinear interaction in the log-likelihood does not imply a rank-one likelihood tensor.
Exact low-rank contraction needs a property of the actual likelihood table, not just of the Gram interaction.

\textbf{What has been gained.} We now have a precise safe state-reduction rule, a direct connection to polar
characters and resource elimination, and an obstruction to an unjustified exact low-rank shortcut.
The remaining question is the number of reachable projected states at useful cuts in this particular bank.
Neither full-bank observation rank nor this counterexample answers that more local question.

\clearpage
\section{Turn the exact cut reduction into a compatible schedule}
\label{sec:nested-schedule}
\subsection{Future spaces shrink along each legal branch}
The preceding projection rule describes one cut. To obtain an algorithm, we must also show how to move between
cuts without losing information or multiplying the same likelihood twice.
Order the patterns and eliminate one entire pattern at a time. A local choice is a payload-count vector $a_p$,
with occupancy $m=\one^Ta_p$, waveform $x=\sum_v a_{p,v}\phi(p,v)$ and weight
$\beta_p(a_p)=\prod_v1/a_{p,v}!$ for IID sampling. For distinct sampling, require binary local counts and use
weight one. Patterns form disjoint message alphabets, so these local prior weights multiply exactly.

After $i$ patterns, let $q$ be the number of assigned users and $h=K-q$ the remaining number.
Let $t_{i,h}$ be any point in the affine hull of the legal remaining sums, and let $P_{i,h}$ project onto their
variation space. These are the objects already derived at a fixed cut. For a nonempty remaining alphabet $S_i$,
one may take $t_{i,h}=h|S_i|^{-1}\sum_{w\in S_i}\phi(w)$. Handle empty or singleton-completion cases as above.

Consider a legal transition from the old cut to the new cut by choosing $a_p$. Write $P_o,t_o$ for the old
geometry and $P_n,t_n$ for the new geometry, whose remaining occupancy is $h-m$.
Every new future sum $t$ gives a legal old future sum $x+t$. Therefore
\begin{equation}
 \operatorname{range}P_n\subseteq\operatorname{range}P_o,\qquad
 x+t_n-t_o\in\operatorname{range}P_o.
 \label{eq:nested-future-spaces}
\end{equation}
The first statement follows by subtracting two such completions. The second follows because affine averages
of legal old completions remain in their affine hull. Both hold at the distinct-source boundary occupancies too.

For an old partial sum $s$, retain $\theta=P_os$. Define the new state and newly resolved residual by
\begin{equation}
 \theta'=P_n(\theta+x),\qquad
 d=(P_o-P_n)(y-\theta-x-t_n).
 \label{eq:nested-transition}
\end{equation}
Since the spaces are nested, $P_o-P_n$ is itself an orthogonal projector. It selects the directions that could
previously vary but can no longer vary after this local choice.
The previously absorbed likelihood factors obey
\begin{equation}
 \boxed{\eta_n(s+x)=\eta_o(s)\exp[-\|d\|^2/(2\sigma^2)].}
 \label{eq:likelihood-telescope}
\end{equation}
To prove it, split $I-P_n=(I-P_o)+(P_o-P_n)$ into orthogonal projectors.
By \eqref{eq:nested-future-spaces}, the first projected residual is exactly the old one.
The second depends on $s$ only through $\theta$. This proves both the update and its independence from any
discarded part of the old state.

\clearpage
\subsection{A complete projected-state recurrence}
Let $F_i(q,\theta)$ be the sum of past prior weights multiplied by the likelihood directions already resolved
at cut $i$, aggregated at $\theta=P_{i,K-q}s$. Initialize
\[
 F_0(0,0)=\exp\!\left[-\frac{\|(I-P_{0,K})(y-t_{0,K})\|^2}{2\sigma^2}\right].
\]
For every legal local count vector, apply
\begin{equation}
 F_{i+1}(q+m,\theta')\mathrel{+}=
 F_i(q,\theta)\,\beta_p(a_p)\,
              \exp[-\|d\|^2/(2\sigma^2)].
 \label{eq:projected-forward}
\end{equation}
Keep the cardinality $q$ in the key: it determines both future feasibility and the appropriate affine geometry.
Sum over all pattern occupancies, including zero. No transmitted pattern is supplied to this algorithm.
At the final cut only $q=K$ is allowed, the future variation is zero, and $\theta'=0$.
The final value is exactly
\[
 F_P(K,0)=\sum_a\beta(a)g(y-\Phi a)=\mathcal Z.
\]
Indeed the incremental orthogonal factors in \eqref{eq:likelihood-telescope} telescope to the single physical
Gaussian. This supplies an exact elimination schedule, not a sequence of independent pseudo-observations.

\subsection{Recover all count queries through the same transitions}
The scalar forward pass alone gives evidence. Let $T_i(s,s')$ denote its total transition weight between
states $s=(q,\theta)$ and $s'$, summing local choices that induce the same transition.
Initialize the backward value $B_P(K,0)=1$ and recurse
\[
 B_i(s)=\sum_{s'}T_i(s,s')B_{i+1}(s').
\]
For a message $w$ in the pattern eliminated at step $i$, define $T_i^{w,j}$ by retaining only local choices
with $a_w=j$. Then
\begin{equation}
 \boxed{\Pr(a_w=j\mid y)=
 \frac{1}{\mathcal Z}\sum_{s,s'}F_i(s)\,T_i^{w,j}(s,s')\,B_{i+1}(s').}
 \label{eq:frontier-count-query}
\end{equation}
Every full count configuration has exactly one such transition, so it contributes once, with its original weight.
An equivalent implementation carries a marker $z^{a_w}$ through the forward recurrence.
If several local multisets share a waveform, their scalar weights may merge, but their marked weights or
$T_i^{w,j}$ contributions must remain available. This is the same identity-preservation requirement as before.

\textbf{What is complete here.} The cut reduction now has a sequential update, initialization, terminal evidence,
and every count query needed by the Bayes decision. It remains potentially exponential. Computing an exact
projector or testing equality of real-valued projected states is itself a representation problem; a numerical
rank threshold would introduce an additional approximation.

\clearpage
\subsection{A simpler executable specialization: finalize each resource once}
Instead of the full affine projection, retain each resource touched by both a processed pattern and an
unprocessed pattern. This set is the \emph{resource frontier}. It is a conservative state: it may retain more
coordinates than necessary, but it loses no information.

When a pattern is processed, add its local waveform to the stored partial sums. If this is the last pattern
that can touch resource $r$, multiply the transition by
$\exp[-(y_r-s_r)^2/(2\sigma^2)]$ and discard that coordinate. Every physical resource is finalized exactly once.
Resources no pattern ever touches contribute their noise-only likelihood at initialization.
The same forward/backward transition formula gives all count laws.
Knowing the remaining occupancy can allow additional reductions; this simple specialization uses support
frontiers to keep the implementation transparent.

The finite reference in \path{tests/polar_exact_computation.py} implements this procedure, including local
equal-signal aggregation and both source laws. It represents the supplied binary floating point amplitudes
as integers with a common power-of-two denominator when constructing sum keys. Thus state merging introduces
no tolerance or extra quantization of the supplied numerical encoder. Gaussian weights still use floating
arithmetic, so the implementation is not a formal interval certificate for the ideal real-valued construction.

For the existing balanced $P=4,V=4,n=12,K=3$ fixture, all 24 pattern orders agree with exhaustive inference.
The best and worst orders by peak forward state count are:
\begin{center}
\begin{tabular}{@{}llrrr@{}}\toprule
Source & Pattern order & Peak forward states & Stored transitions & Peak frontier\\\midrule
IID & $(0,3,2,1)$ & 286 & 1,202 & 8\\
IID & $(1,2,3,0)$ & 455 & 1,471 & 8\\
Distinct & $(0,3,2,1)$ & 196 & 788 & 8\\
Distinct & $(1,2,3,0)$ & 299 & 967 & 8\\\bottomrule
\end{tabular}
\end{center}
This fixture has 816 IID or 560 distinct full count configurations.
The table reports reachable states and stored transitions, not wall-clock speedup or total-memory equivalence.
The forward/backward implementation stores the transition graph as well as forward values.
Even with the same largest frontier dimension, the number of reachable values changes substantially with order.

These calculations establish that compatible elimination and lossless merging can be executed and checked.
They do not establish useful scaling at the selected $B=12/14$ sizes. The next section examines those banks
structurally, without attempting to enumerate their multiuser states.

\clearpage
\section{How many mixed-pattern states can remain distinguishable?}
\label{sec:mixed-state-counts}
\subsection{Private coordinates give a lower bound without enumeration}
Consider a cut whose processed pattern set is $A$, with exactly one user on each of its $q$ patterns.
This is a subfamily of legal partial hypotheses under both source laws. Suppose coordinate $r$ is used by
pattern $p\in A$ but by no other pattern in $A$. Its partial signal is
$\sqrt{e_{p,i}}\,(-1)^{\ell_i^Tv_p}$, so a positive known amplitude leaves its parity bit identifiable.
Use only such coordinates whose coordinate axes are known to lie in the future variation space.

Let $r_p$ be the binary rank of the corresponding rows of $L$. Those coordinates have $2^{r_p}$ different
sign patterns as $v_p$ varies. The private coordinate sets are disjoint across $p$, so their combinations give
\begin{equation}
 \boxed{\text{number of distinct retained partial states}\ \geq\
             2^{\sum_{p\in A}r_p}.}
 \label{eq:private-state-bound}
\end{equation}
Any change in one of these coordinates remains visible after future-space projection, since its axis belongs
to that space. The proportional-likelihood criterion then prevents merging the resulting states up to a scalar.
Other partial occupancies may create additional states; this lower bound needs only the specified subfamily.

\subsection{Apply the certificate to the actual selected banks}
For all six selected banks, after processing the first eight patterns, the remaining pattern alphabet still
has affine variation $\R^{256}$. The support-based axis certificate proves this, independently of numerical
rank thresholds. Each of the eight processed patterns has private-row rank six. Hence there are at least
\[
 64^8=2^{48}=281{,}474{,}976{,}710{,}656
\]
distinguishable partial states at this cut. This applies when future occupancy is positive: IID $K>8$, or
distinct sampling with $1\leq K-8<64(P-8)$. At zero future occupancy all signal directions may instead be
absorbed into the final likelihood; that is a different cut.

The same calculation at other natural-order cuts gives the following exponents in \eqref{eq:private-state-bound}:
\begin{center}
\begin{tabular}{@{}lrrrr@{}}\toprule
Seed, either $B=12$ or $14$ & $q=4$ & $q=8$ & $q=12$ & $q=16$\\\midrule
3201 & 24 & 48 & 65 & 62\\
3202 & 24 & 48 & 68 & 67\\
3203 & 24 & 48 & 64 & 65\\\bottomrule
\end{tabular}
\end{center}
At all these cuts the future axis certificate gives all 256 axes. The exponents need not increase with $q$:
adding patterns removes some previously private coordinates. They are certified lower bounds, not exact counts.

\textbf{Scope of this obstruction.} An explicit table that only merges proportional continuation rows is enormous
at these cuts. A partial hypothesis still has a compact description, and a factored or symbolic calculation may
avoid materializing the table. This is not a proof that every exact decoder needs this much memory, nor a
complexity lower bound for all elimination orders.

\clearpage
\subsection{Resource-frontier ordering helps, but does not make these banks narrow}
The exact resource-frontier widths can be counted without decoding.
Compare natural pattern order with a deterministic greedy order choosing the next pattern that minimizes
the resulting frontier width, breaking ties by pattern index:
\begin{center}
\begin{tabular}{@{}lrrr@{}}\toprule
$B$ & Seed & Natural peak width & Greedy peak width\\\midrule
12 & 3201 & 248 & 234\\
12 & 3202 & 251 & 234\\
12 & 3203 & 252 & 231\\
14 & 3201 & 256 & 254\\
14 & 3202 & 256 & 254\\
14 & 3203 & 256 & 254\\\bottomrule
\end{tabular}
\end{center}
The greedy rule chooses an evaluation order; it does not approximate inference.
It is not claimed to minimize frontier width globally. These two schedules leave wide frontiers, while the
tiny calculation already showed that width alone does not determine the number of reachable states.

\subsection{A precise minimum for linear contraction, with a narrower scope}
Fix finite partial and completion signal sets $\{s_i\}$ and $\{t_j\}$ at fixed cardinalities. Their coupling is
the matrix $G_{ij}=g(y-s_i-t_j)$, and evidence contracts it as $a^TGb$.
If a linear state $a^TU$ followed by a linear readout $Vb$ must preserve this for \emph{all} partial and
completion weights, then $G=UV$. The smallest possible intermediate dimension is exactly
\[
 r=\operatorname{rank}G.
\]
The lower bound follows from $\operatorname{rank}(UV)\leq r$; a rank factorization attains it.
Polynomial count markers can pass through the same linear factorization coefficient by coefficient.
Allowing signed factors may introduce cancellation; requiring nonnegative factors is an additional restriction.

Observation-dependent row and column scalings do not change this rank:
\[
 G_{ij}=e^{-\|y\|^2/(2\sigma^2)}
 e^{(y^Ts_i-\|s_i\|^2/2)/\sigma^2}
 \underbrace{e^{-s_i^Tt_j/\sigma^2}}_{E_{ij}}
 e^{(y^Tt_j-\|t_j\|^2/2)/\sigma^2}.
\]
Thus exact cut rank is independent of $y$ at fixed signal sets and noise variance.
Disjoint signal subspaces give a rank-one coupling; the scalar Vandermonde example gives arbitrarily large rank.
Numerically small singular values would justify an approximation only after an error analysis; they are not
exact zeros.

This separates three questions: equality of continuation rows, linear dependence of different rows, and the cost
of evaluating a structured representation without enumerating its rows. The $2^{48}$ certificate settles the
first at one concrete cut. It does not settle the second or third for the full bank.

\clearpage
\section{Contract polar interactions through their overlap parities}
\label{sec:pair-contraction}
\subsection{What one pattern pair actually needs to know}
The large-state obstruction motivates a more local question: can a required matrix operation be performed
without forming its full message-pair matrix? For the polar encoder, the answer is yes.
This is an exact structured contraction, even when the matrix is full rank.

Let $I_{pq}$ be the resources shared by patterns $p$ and $q$, and let $i_p(r)$ be the local symbol placed
on shared resource $r$. Their Gram entry is
\begin{equation}
 c_{pq}(v,u)=\phi(p,v)^T\phi(q,u)
 =\sum_{r\in I_{pq}}\sqrt{e_{p,i_p(r)}e_{q,i_q(r)}}\,
       (-1)^{\ell_{i_p(r)}^Tv+\ell_{i_q(r)}^Tu}.
 \label{eq:overlap-gram}
\end{equation}
Thus it depends on $v$ only through the parity vector
$\alpha_{pq}(v)=(\ell_{i_p(r)}^Tv)_{r\in I_{pq}}$, and on $u$ only through $\alpha_{qp}(u)$.
All parity arithmetic here is binary. Put
\[
 r_{pq}=\operatorname{rank}_{\mathrm{GF}(2)}L_{I_{pq},p}.
\]
There are exactly $2^{r_{pq}}$ attainable left parity vectors, regardless of how many repeated or constant
generator rows occur on the overlap. Enumerating just these vectors yields a table with
$2^{r_{pq}+r_{qp}}$ entries instead of $V^2$, where $V=2^{B_c}$.

The same grouping applies to the exact pair factor $b=e^{-c/\sigma^2}$, its signed correction $f=b-1$,
and $F=|f|$. For any one of these kernels $J_{pq}$ and any input weights $x_q(u)$, first form
\[
 X_q(\beta)=\sum_{u:\alpha_{qp}(u)=\beta}x_q(u),
 \qquad
 (J_{pq}x_q)(v)=\sum_\beta\overline J_{pq}(\alpha_{pq}(v),\beta)X_q(\beta).
\]
This is an exact identity for arbitrary signed weights. It requires grouping, multiplication by the smaller
table, and copying the result to payloads with the same left parity vector. No singular-value threshold is used.
The table is an upper bound on required representation size, not a claim of minimum rank.

\subsection{Within one pattern, use XOR convolution}
For $p=q$, the Gram entry is a function only of payload XOR:
\[
 c_{pp}(v,u)=\sum_i e_{p,i}(-1)^{\ell_i^T(v\oplus u)}.
\]
Consequently $J_{pp}(v,u)=\kappa_p(v\oplus u)$. If
$\widehat x(g)=\sum_v(-1)^{g^Tv}x(v)$ denotes the unnormalized Walsh transform, then
\begin{equation}
 J_{pp}x=\frac1V\operatorname{Walsh}(\widehat\kappa_p\,\widehat x).
 \label{eq:pair-walsh-contraction}
\end{equation}
Applying the unnormalized Walsh operation twice gives $V$ times the input; the factor $1/V$ appears once.
This costs $O(V\log V)$ arithmetic and stores $V$ kernel values.
For distinct sampling, set $b_{pp}(v,v)=0$, hence $f_{pp}(v,v)=-1$ and $F_{pp}(v,v)=1$.
Only the zero-XOR kernel value changes, so the convolution remains valid.

\clearpage
\subsection{Assemble the full unknown-pattern operator}
Sum the pattern-pair contributions over every $q$ for each output pattern $p$. This retains all unknown
patterns. Disjoint patterns have $c=0$, so their $f$ and $F$ blocks vanish. For $b$, add the all-ones
background: $bx=\one(\one^Tx)+fx$.
The finite reference implements both $f$ and $F$ in \path{tests/polar_exact_reuse.py}.
Signed corrections and Walsh arithmetic can suffer cancellation; exact algebra is not a controlled-rounding
machine certificate.

The selected-bank census counts all unordered intersecting pairs. The dense comparison already exploits
symmetry and omits disjoint blocks:
\begin{center}
\begin{tabular}{@{}lrrr@{}}\toprule
$B$ & Seed & Overlap-parity table entries & Dense/table ratio\\\midrule
12 & 3201 & 1,130,735 & 7.24\\
12 & 3202 & 1,149,633 & 7.13\\
12 & 3203 & 1,176,303 & 6.96\\
14 & 3201 & 19,375,511 & 6.83\\
14 & 3202 & 19,318,283 & 6.85\\
14 & 3203 & 19,133,478 & 6.92\\\bottomrule
\end{tabular}
\end{center}
In addition, same-pattern kernels need $PV$ values and parity-group maps require indices.
The ratio concerns pair-table entries, not total memory or measured speedup. Each cross-pattern multiply
also groups and copies $V$ weights at each endpoint. This remains explicit in the pattern and payload alphabets.

These are local contractions. Eliminating a variable with several neighbours can create a larger joint factor;
the small pair tables do not make that joint factor disappear. Their concrete use below is to evaluate weighted
tree bounds without building the dense $M\times M$ interaction matrix.

\subsection{A complete exact two-user decoder follows immediately}
For $K=2$, this operator also completes inference itself. Use the one-user reference
$\nu_w\propto e^{y^T\phi(w)/\sigma^2}$ and the exact pair factor $b$ for the chosen source law.
Compute
\[
 s=b\nu,\qquad Z_2=\nu^Ts.
\]
The ordered two-slot posterior is $\nu_u\nu_vb(u,v)/Z_2$.
Both slot permutations contribute when a word occurs once; there is only one ordered pair when both slots
contain that word. Therefore
\begin{equation}
 q_{w,2}=\frac{\nu_w^2b(w,w)}{Z_2},\qquad
 q_{w,1}=\frac{2\nu_ws_w}{Z_2}-2q_{w,2},\qquad
 q_{w,0}=1-q_{w,1}-q_{w,2}.
 \label{eq:structured-two-user}
\end{equation}
For IID, $b(w,w)=e^{-1/\sigma^2}$; for distinct sampling it is zero.
Use these probabilities to form the count tails and allocate two output copies as in Section~3.
This includes unknown patterns and same-pattern collisions. The finite test compares every count probability
with exhaustive inference. The next section handles the first cycle, at $K=3$.

\clearpage
\section{Three users: close the triangle exactly}
\label{sec:three-user}
\subsection{Three pair interactions, with no new physical assumptions}
Keep the fixed encoder, known $K=3$, Gaussian likelihood and either source law. Write
\[
 \nu_w=\frac{e^{y^T\phi(w)/\sigma^2}}{\sum_u e^{y^T\phi(u)/\sigma^2}},\qquad
 B_{uv}=e^{-\phi(u)^T\phi(v)/\sigma^2}
 \begin{cases}1,&\text{IID},\\ \ind\{u\ne v\},&\text{distinct}.\end{cases}
\]
Here $B$ is the $M\times M$ pair matrix, not the message-bit length; subscripts distinguish its entries.
Expanding the squared residual as in \eqref{eq:interactions} gives
\begin{equation}
 \Pr(w_1=u,w_2=v,w_3=t\mid y)
 =\frac{\nu_u\nu_v\nu_t B_{uv}B_{vt}B_{tu}}{Z_3}.
 \label{eq:triangle-posterior}
\end{equation}
Every pair of slots interacts, so the graph is a triangle. The normalized one-user weights $\nu$ are only a
convenient reference; they are not the actual marginals. Constant likelihood and source-normalization factors
cancel. All pattern choices and allowed full-message repeats remain in this sum.

Put $D_\nu=\diag(\nu)$ and $A=D_\nu^{1/2}BD_\nu^{1/2}$. Then
\begin{equation}
 \boxed{Z_3=\tr(A^3),\qquad \E[a_w\mid y]=\frac{3(A^3)_{ww}}{Z_3}.}
 \label{eq:triangle-trace}
\end{equation}
Indeed, $A_{uv}A_{vt}A_{tu}=\nu_u\nu_v\nu_t B_{uv}B_{vt}B_{tu}$; the trace sums every ordered triangle.
Fixing its first vertex gives the diagonal entry, and exchangeability supplies the factor three.
$Z_3$ is the normalizer of this reference-weight expression, not the physical evidence including its prefactor.
The identity does not require $A$ to be positive semidefinite, and it does not justify discarding eigenmodes.
For IID, expected counts alone are still insufficient for minimum multiset-PUPE decoding.

\subsection{Where the first joint factor appears}
Eliminating the middle slot leaves the endpoint factor
\begin{equation}
 C(u,t)=\sum_v\nu_v B_{uv}B_{vt}=(BD_\nu B)_{ut},\qquad
 \Pr(u,t\mid y)=\frac{\nu_u\nu_t B_{ut}C(u,t)}{Z_3}.
 \label{eq:triangle-separator}
\end{equation}
Both endpoints depend on the \emph{same} unknown middle message. In particular,
\[
 C(u,t)=(B\nu)_u(B\nu)_t+
 \operatorname{Cov}_{V\sim\nu}(B_{uV},B_{tV}).
\]
Replacing the joint factor by a product drops this covariance. A fast pair operator does not by itself make
the covariance vanish or make the joint factor small. We can nevertheless obtain the needed count queries
without storing the entire joint table.

\clearpage
\subsection{A complete decoder from one streamed contraction per message}
For a specified full message $w$, let $\lambda_w=B_{ww}$ and define
\[
 x_w(u)=\nu_u B_{wu}\ind\{u\ne w\}.
\]
Separate the ordered triples by whether they contain one, two or three copies of $w$. Their unnormalized
posterior masses are
\begin{align}
 N_1(w)&=3\nu_w\,x_w^TBx_w,\label{eq:triangle-one}\\
 N_2(w)&=3\nu_w^2\lambda_w\sum_{u\ne w}\nu_uB_{wu}^2,
 \label{eq:triangle-two}\\
 N_3(w)&=\nu_w^3\lambda_w^3.\label{eq:triangle-three}
\end{align}
For $N_1$, choose which of the three slots contains $w$ and sum the other two slots, both excluding $w$.
Their mutual interaction remains in $B$. For $N_2$, choose the singleton slot: the repeated pair contributes
$\lambda_w$, and its two interactions with the singleton give $B_{wu}^2$. For $N_3$, all three edges are
diagonal. These are positive sums, including same-pattern collisions under IID.

Every triple has exactly three message occurrences, so its weight is counted three times in
$\sum_w\sum_{j=1}^3 jN_j(w)$. Therefore
\begin{equation}
 Z_3=\frac13\sum_w\{N_1(w)+2N_2(w)+3N_3(w)\},\qquad
 q_{w,j}=\frac{N_j(w)}{Z_3}\quad(j=1,2,3),\quad
 q_{w,0}=1-\sum_{j=1}^3q_{w,j}.
 \label{eq:triangle-counts}
\end{equation}
For distinct sampling $\lambda_w=0$, automatically giving $q_{w,2}=q_{w,3}=0$. For IID,
$\lambda_w=e^{-1/\sigma^2}$ and all three terms are required.
Form the tails $T_{w,j}=\sum_{r=j}^3q_{w,r}$ and allocate three copies by Section~3. This completes the exact
Bayes decoder, including the output decision, rather than only evaluating its evidence.

For each $w$, generate one row of $B$, construct $x_w$, apply the structured pair operator to it, and save
the three numerators. Normalize after processing every $w$. This uses $M$ pair-operator applications,
$O(M^2)$ additional scalar work, and $O(M)$ working storage \emph{in addition to} the operator's stored
kernels, overlap tables and maps. No $M^3$ tensor or full $M^2$ pair table is required. The count vector still
has $M=2^B$ entries, and the repeated contractions may be expensive. This is a memory reduction and an exact
computational specification, not a full-bank runtime result.

\textbf{Numerical representation matters.} Algebraically, $Bx=\one\sum_u x_u+fx$. Numerically, this can subtract
nearly equal quantities. At $\sigma^2=.02$, the IID unit-energy diagonal is $e^{-50}>0$, but ordinary double
precision evaluates $1+\operatorname{expm1}(-50)$ as zero. The new positive reference uses direct $B$ overlap
tables and positive XOR sums within a pattern, costing $O(PV^2)$ there instead of the faster Walsh arithmetic.
No cancellation-prone reconstruction from $f$ is needed. Extreme underflow or overflow still requires explicit
rescaling or higher precision; the checks below are floating-point comparisons, not rounding certificates.

\clearpage
\subsection{Combine both neighbours' parity requirements before grouping}
There is also a direct polar-specific contraction for a fixed ordered pattern triple $(p_1,p_2,p_3)$.
For slot $i$, let $J_i$ contain the local symbols whose resources overlap \emph{either} other pattern. Define
\[
 \alpha_i(v)=(L_{J_i}v)\bmod2,\qquad
 r_i=\operatorname{rank}_{\GF(2)}L_{J_i},\qquad
 m_i(\alpha)=\sum_{v:\alpha_i(v)=\alpha}\nu_{(p_i,v)}.
\]
Both adjacent pair factors are determined by this one joint signature. The fixed-pattern contribution is
\begin{equation}
 Z_{p_1p_2p_3}=
 \sum_{\alpha_1,\alpha_2,\alpha_3}
 m_1(\alpha_1)m_2(\alpha_2)m_3(\alpha_3)
 \overline B_{12}(\alpha_1,\alpha_2)
 \overline B_{13}(\alpha_1,\alpha_3)
 \overline B_{23}(\alpha_2,\alpha_3).
 \label{eq:triangle-quotient}
\end{equation}
There are $2^{r_1+r_2+r_3}$ attainable signature triples. This counts a direct sum's states, not a lower bound
on every contraction algorithm. Grouping separately for each edge and then treating those groups as independent
would lose the fact that both signatures come from the same payload.

To recover a specified message's count law, replace each unary by the polynomial
\[
 m_i(\alpha;z)=m_i(\alpha)+
 \ind\{p_i=p_w,\ \alpha=\alpha_i(v_w)\}\nu_w(z-1).
\]
Only the weight of message $w$ is multiplied by $z$. The coefficient of $z^j$ in the triangle contraction is
its unnormalized count-$j$ mass. Thus unmarked payloads can share a class while the query retains the marked
message's identity. If two patterns coincide, their shared resources expose the full generator. For the
injective payload code used here, these classes identify the whole payload, so distinct-message exclusions
remain exact. A noninjective generator would need an additional identity state.

Sum over all ordered pattern triples; no true pattern is supplied to the receiver. Because evidence and count
queries are slot-permutation invariant, one may instead sum pattern multisets, weighting each by
$3!/\prod_p n_p!$, where $n_p$ is its occupancy. Same-pattern triples remain included.

\textbf{How much grouping survives?} For the selected $B_c=6,N_c=32$, CRC8 code, enumerate the 56 triples of
\emph{distinct} patterns among the first eight patterns of each $B=12$ bank:
\begin{center}
\begin{tabular}{@{}rrrrr@{}}\toprule
Seed & Minimum states & Median states & Maximum states & All three ranks six\\\midrule
3201 & 1,024 & 131,072 & 262,144 & 14/56\\
3202 & 16,384 & 131,072 & 262,144 & 15/56\\
3203 & 1,024 & 65,536 & 262,144 & 12/56\\\bottomrule
\end{tabular}
\end{center}
An ungrouped payload triple has $64^3=262,144$ states. The joint signature frequently retains much of this
alphabet, and sometimes all of it: for seed 3201, patterns $(0,1,6)$ have $(r_1,r_2,r_3)=(6,6,6)$.
These are exact counts for the stated 168 triples, not a random sample or a census of the full bank.
They illustrate why two small edge interfaces need not compose into an equally small joint interface.

\clearpage
\subsection{An exact identity for what deleting the closing edge loses}
The completed triangle lets us examine a possible future assumption precisely. Define a reference
$Q_{\rm open}$ by removing the \emph{Gaussian} interaction between slots 1 and 3 but keeping their original
unary weights and \emph{all source-law constraints}. Let
\[
 R(W)=e^{-\phi(w_1)^T\phi(w_3)/\sigma^2}.
\]
For IID, the reference has a chain factor graph. For distinct sampling it still enforces $w_1\ne w_3$ as
well as the two other exclusions; it is not a tree model with independent endpoints. On either legal support,
\begin{equation}
 P(W\mid y)=\frac{Q_{\rm open}(W)R(W)}{\E_{\rm open}R},\qquad
 \E_P H-\E_{\rm open}H
 =\frac{\operatorname{Cov}_{\rm open}(H,R)}{\E_{\rm open}R}.
 \label{eq:triangle-closure-covariance}
\end{equation}
The second identity follows by expanding the covariance. With
$H=\ind\{a_w\geq j\}$, it gives the exact error in every count tail. It does not assume the deleted edge is
weak. Small unweighted Gram correlations alone do not bound this observation-weighted covariance.

The reference singles out slot endpoints and need not preserve exchangeability. Count observables must
therefore aggregate all three slots. The covariance identity is retained as an exact comparison resource;
no edge-deletion approximation or empirical regret ranking is adopted in this report.

\subsection{Why the trace shortcut stops at three users}
The three-user formula should not be extrapolated as $Z_K=\tr(A^K)$. A trace of a matrix power closes a
\emph{cycle}; the physical Gaussian model contains every pair, a \emph{complete graph}. These coincide for
three vertices. For four, the cycle lacks two interactions. Under distinct sampling it also misses the
corresponding full-message exclusions.

A simple exact counterexample uses four orthogonal unit waveforms, $y=0$ and distinct sampling. Then
$\nu_w=1/4$, $B=\one\one^T-I$, and
\[
 Z_4=\frac{4!}{4^4}=\frac{24}{256},\qquad
 \tr\big((B/4)^4\big)=\frac{3^4+3}{4^4}=\frac{84}{256}.
\]
The cycle sum includes repeats at opposite vertices; the physical distinct-source posterior forbids them.
Higher $K$ can still be solved by the exact constructions already given, but closing a single chain is no
longer sufficient. This is a boundary of this shortcut, not an impossibility theorem for other reductions.

The general finite decoder, the complete first-cycle contraction and the exact growth of joint parity
requirements are now explicit. Further assumptions can refer to these calculations and the covariance
identities without treating pairwise interfaces as a complete solution for larger K. Report 9 owns that
algebraic investigation; the present report continues with exact decision certificates.

\clearpage
\section{Certify the Bayes decision without finishing every sum}
\label{sec:certificate}
\subsection{Positive bounds on unexplored count subtrees}
Another exact endpoint is a finite search that stops when its \emph{decision} is proved.
Order the $M$ full messages and build a tree assigning their counts.
A node has fixed counts $c_1,\ldots,c_d$, remaining users $h=K-\sum_{i\leq d}c_i$,
remaining alphabet size $m=M-d$, and residual $r=y-\sum_{i\leq d}c_i\phi_i$.
Its completions have total signal norm at most $h$, since each message has unit norm. Hence
\[
 \|r-s_{\rm remaining}\|\geq \max(0,\|r\|-h).
\]
The sum of the coefficients $\beta(a)$ over the node's completions is known exactly:
\begin{equation}
 B_{\rm node}=
 \begin{cases}
 \displaystyle\frac{m^h}{h!\prod_{i\leq d}c_i!},&\text{IID},\\[1mm]
 \displaystyle\binom mh,&\text{distinct and the prefix is valid}.
 \end{cases}
 \label{eq:subtree-prior}
\end{equation}
For IID this is the multinomial theorem; for distinct sampling it counts subsets.
Use $0^0=1$ and zero weight for infeasible nodes. The remaining evidence mass is bounded by
\begin{equation}
 U_{\rm node}=B_{\rm node}
 \exp\!\left[-\frac{\max(0,\|r\|-h)^2}{2\sigma^2}\right].
 \label{eq:subtree-bound}
\end{equation}
This ball bound is often loose. Its advantage is a short proof and no oracle assumptions.
It remains valid when the unexplored messages overlap in resources or share a pattern.

Let $\mathcal E$ contain evaluated leaves, and let $\mathcal F$ be a disjoint frontier covering every other
valid state exactly once. With $\omega(a)=\beta(a)g(y-\Phi a)$, define
\[
 Z_-=\sum_{a\in\mathcal E}\omega(a),\qquad
 Z_+=Z_-+\sum_{b\in\mathcal F}U_b.
\]
Then $Z_-\leq\mathcal Z\leq Z_+$. Refining a frontier node replaces its bound by bounds on its children or by
exact leaf weights. No posterior state is silently discarded.

\subsection{Bound only what the loss asks for}
Define the unnormalized count-tail numerator
$N_{w,j}=\sum_a\omega(a)\ind\{a_w\geq j\}$, so $T_{w,j}=N_{w,j}/\mathcal Z$.
Simple valid bounds are
\begin{align}
 L_{w,j}&=\sum_{a\in\mathcal E}\omega(a)\ind\{a_w\geq j\},\\
 U_{w,j}&=L_{w,j}
     +\sum_{\substack{b\in\mathcal F:\\a_w\geq j\text{ possible in }b}}U_b .
 \label{eq:tail-bounds}
\end{align}
Possibility follows from the fixed prefix and remaining count budget, without decoding a completion.
Because every candidate has the same positive denominator $\mathcal Z$, ranking these unnormalized
numerators is sufficient. Computing accurate evidence normalization is not required to certify the output.

\clearpage
\subsection{A sufficient stopping rule, and why it proves optimality}
For a proposed allocation $\widehat a$ with $\sum_w\widehat a_w=K$, test
\begin{equation}
 \boxed{\min_{w:\widehat a_w>0}L_{w,\widehat a_w}
 \ \geq\
 \max_{w:\widehat a_w<K}U_{w,\widehat a_w+1}.}
 \label{eq:decision-certificate}
\end{equation}
The left side bounds each message's last selected occurrence from below; the right side bounds each message's
next unselected occurrence from above. Earlier tails of a message are no smaller than its last selected tail,
and later tails are no larger than its next unselected tail. Therefore every selected occurrence has at least
as much posterior value as every unselected occurrence. This proves a Bayes-optimal allocation, including ties.
For the distinct model, tails beyond the first are zero.

An exact search can now repeat: evaluate or split an unresolved node, update the bounds, propose an allocation
by greedy lower tails, and test \eqref{eq:decision-certificate}. If certification fails, continue.
After all leaves are evaluated, lower and upper bounds coincide and the ordinary Bayes allocation certifies.
Thus exhaustive completion is a guaranteed finite fallback in the mathematical algorithm, not a heuristic
candidate list. Some observations can certify earlier; no typical-case speedup has been established.

For a concrete illustration, take $K=2$ and unnormalized tail intervals
\[
 N_{u,1}\in[8,9],\quad N_{u,2}\in[1,2],\qquad
 N_{v,1}\in[6,7],\quad N_{v,2}\in[0,1],
\]
with every other first-tail upper bound at most $3$. Returning one $u$ and one $v$ is certified:
the smallest selected lower bound is $6$ and the largest unselected upper bound is at most $3$.
Neither the exact numerators nor their common normalizer need be known.
These intervals are illustrative, not a measured decoder outcome.

\subsection{A loss bound when the exact decision is not yet certified}
Write $A(\widehat a)=\{(w,j):1\leq j\leq\widehat a_w\}$. If $Z_->0$, a conservative certificate is
\begin{equation}
 0\leq R(\widehat a\mid y)-R^*(y)
 \leq \min\!\left\{1,\,
 \frac{\displaystyle\max_{A:\,|A|=K,\ {\rm prefixes}}\sum_{(w,j)\in A}U_{w,j}
       -\displaystyle\sum_{(w,j)\in A(\widehat a)}L_{w,j}}
      {K Z_-}\right\}.
 \label{eq:risk-certificate}
\end{equation}
The proof replaces the unknown optimal numerator sum by an upper bound, the proposed sum by a lower bound,
and the positive denominator by $Z_-$. The bounds in \eqref{eq:tail-bounds} are nonincreasing in $j$,
so the maximization over upper tails uses the same prefix-respecting greedy allocation.
This bounds \emph{decision loss}; it need not certify an accurate approximation of the entire posterior.

\textbf{Exactness and the next theoretical question.} A complete exact decoder has now been derived twice:
positive convolution with outside queries, or exhaustive refinement until a decision certificate holds.
Polar structure provides lossless signatures and a one-user transform, but no affordable general separator
state or sharp typical-case complexity bound yet. Tightening those exact states or bounds is still a valid
next step. Neither the algebra nor this stopping point forces a switch to any of the three routes below.
If learning is considered later, the evidence functions or decision bounds identify concrete quantities to
approximate; no learned component is required here.

\clearpage
\section{Tighter completion bounds from the encoding rule}
\label{sec:completion-dual}
\subsection{Use the convex hull only to bound unfinished inference}
Return to a count-tree node with residual $r$, remaining alphabet $S$ and remaining count $h$.
Let $\mathcal T_{S,h}$ contain its possible completion signals. The exact minimum residual is
$d_{\rm int}^2=\min_{t\in\mathcal T_{S,h}}\|r-t\|^2$.
Replacing this finite set by its convex hull can only decrease that minimum. Write
$\mathcal C_{S,h}=\operatorname{conv}\mathcal T_{S,h}$ and define its support function
\[
 H_{S,h}(u)=\max_{t\in\mathcal C_{S,h}}u^Tt.
\]
It asks a simple question: how far can a legal completion point in direction $u$?
Its exact form follows from the source law:
\begin{equation}
 H_{S,h}(u)=
 \begin{cases}
 h\max_{w\in S}u^T\phi(w),&\text{IID},\\
 \displaystyle\sum_{\ell=1}^h s_{(\ell)}(u),
       \quad s_w(u)=u^T\phi(w),&\text{distinct},
 \end{cases}
 \label{eq:completion-support}
\end{equation}
where $s_{(1)}\geq s_{(2)}\geq\cdots$ are scores of distinct messages in $S$.
Set $H_{S,0}=0$; discard infeasible nodes. IID can repeat the maximizing message. Distinct sampling must choose
the $h$ different maximizing messages. Equivalently, the two hulls are
\[
 \mathcal C_{S,h}^{\rm IID}=\{\Phi_Sx:x\geq0,\ \one^Tx=h\},\qquad
 \mathcal C_{S,h}^{\rm dist}=\{\Phi_Sx:0\leq x\leq1,\ \one^Tx=h\}.
\]
The second polytope is the convex hull of cardinality-$h$ indicator vectors: a fractional vertex would have
at least two fractional coordinates, which can be perturbed in opposite directions while preserving the sum.

For any trial vector $u$, completing a square gives
\[
 \|r-t\|^2
 =\|r-t-u\|^2+2u^T(r-t)-\|u\|^2
 \geq 2u^Tr-\|u\|^2-2H_{S,h}(u).
\]
Therefore
\begin{equation}
 \boxed{L(u)=2u^Tr-\|u\|^2-2H_{S,h}(u)\ \leq\ d_{\rm int}^2.}
 \label{eq:completion-dual-bound}
\end{equation}
No convergence of an optimizer is needed for this inequality. One valid trial vector already provides a bound.
Combining trials with the old ball bound yields
\[
 L_{\rm node}=\max\{0,\ \max(0,\|r\|-h)^2,\ L(u_1),\ldots,L(u_J)\},\qquad
 U_{\rm node}=B_{\rm node}e^{-L_{\rm node}/(2\sigma^2)}.
\]
Use this upper evidence bound in exactly the same frontier, count-tail and decision certificates as before.
The count prior has not been relaxed or changed.

\clearpage
\subsection{What optimizing the bound achieves}
The best dual bound is precisely the convex distance:
\begin{equation}
 \sup_u L(u)=\min_{t\in\mathcal C_{S,h}}\|r-t\|^2.
 \label{eq:completion-dual-equality}
\end{equation}
Here is a direct proof. Let $t_*$ be the Euclidean projection of $r$ onto the compact convex hull.
Optimality along each line segment from $t_*$ implies
$(r-t_*)^T(t-t_*)\leq0$ for every $t$ in the hull. Thus for $u_*=r-t_*$,
$H_{S,h}(u_*)=u_*^Tt_*$. Substitution into $L$ gives $\|r-t_*\|^2$.
Together with the lower-bound inequality, this proves equality.

The convex distance need not equal the finite-set distance. It is a bound used within an exact decoder,
not a replacement of the posterior by continuous counts. A numerical convex solver can suggest $u$;
evaluate a valid support upper bound afterward to establish a valid lower bound on distance.
An approximate primal objective alone has the wrong direction for this purpose.

\subsection{The exact polar interface returns in a second role}
At a fixed pattern, each support score is
$u^TA_pb(v)$. The maximum can therefore use the exact polar maximization operation from Section 4,
with its appropriate payload-prefix or exclusion constraints. Across patterns take the maximum; for the
distinct support, obtain the exact best $h$ different full messages across all patterns.
Exhaustive payload scores or an exact search suffice, even when computationally expensive.

An upper bound $\overline H(u)\geq H_{S,h}(u)$ is also safe: replacing $H$ by $\overline H$ decreases $L$.
For example, searching a superset of $S$ gives a valid, possibly weaker bound.
By contrast, an incomplete candidate list produces a \emph{lower} estimate of the maximum and can make the
certificate false. Ordinary bounded SCL output therefore cannot be inserted as a support oracle without an
additional proof bounding its omitted candidates.

Two small examples show what the stronger geometry can accomplish. Let the available unit waveforms be the
two coordinate vectors $e_1,e_2$, with $h=2$:
\begin{center}
\begin{tabularx}{\textwidth}{@{}lcccY@{}}\toprule
Residual & Ball bound & IID hull & Distinct hull & Explanation\\\midrule
$(-1,-1)$ & $0$ & $8$ & $8$ & Positive completions cannot approach a negative residual.\\
$(2,0)$ & $0$ & $0$ & $2$ & Repetition reaches the residual; distinct selection cannot.\\\bottomrule
\end{tabularx}
\end{center}
Entries are squared-distance lower bounds. For the first row, $u=(-2,-2)$ already attains 8.
For the second distinct row, $u=(1,-1)$ attains 2.
Neither example depends on learning, a posterior approximation or an oracle for the transmitted messages.

\textbf{What remains to investigate.} A useful implementation must balance the cost of improving $u$ and its
support bound against simply splitting the count-tree node. The exact derivation now permits either choice
while retaining the same stopping guarantee. It does not yet establish a tractable search for the full bank.

\clearpage
\section{An exact Fourier representation of the convolution}
\label{sec:fourier-exact}
\subsection{Change the basis of the state, keeping its information}
The signal-sum measures in Section~\ref{sec:tree-exact} can be represented by their Fourier coefficients.
This is an alternative implementation of the same convolution, not another posterior.
Use the integer polar signatures from \eqref{eq:integer-signature}: write each message as
$\phi(w)=C t_w$, where $t_w\in\{-1,0,1\}^d$ and $C$ contains the known real amplitude/insertion coefficients.
For the direct grouped construction $d=P|\mathcal G|$.

A sum of at most $K$ messages lies in $[-K,K]^d\cap\Z^d$. Choose $Q=2K+1$ and regard this set as the
centered representatives of $(\Z/Q\Z)^d$. Every admissible integer sum has a unique residue.
For frequencies $\ell\in\{0,\ldots,Q-1\}^d$, put
\[
 \theta_\ell=2\pi\ell/Q,\qquad
 \rho_\ell=\sum_{w\in\calW}e^{-{\rm i}\theta_\ell^Tt_w}.
\]
The Fourier transform of $\delta_{t_w}$ is its phase $e^{-{\rm i}\theta_\ell^Tt_w}$.
Multiplying phases adds signatures, so convolution becomes multiplication.
For the cardinality-$K$ measure,
\begin{equation}
 \widehat\mu_K(\ell)=
 \begin{cases}
 \rho_\ell^K/K!,&\text{IID},\\
 [u^K]\displaystyle\prod_{w\in\calW}
       (1+u e^{-{\rm i}\theta_\ell^Tt_w}),&\text{distinct}.
 \end{cases}
 \label{eq:fourier-measure}
\end{equation}
$[u^K]$ means the coefficient of $u^K$. The IID formula is the transform of $K$ ordered draws, divided by
$K!$ to give the count weights $1/\prod_w a_w!$. The distinct product chooses each message zero or one time.

\subsection{Invert first, then apply the physical likelihood}
The finite inverse transform is
\begin{equation}
 \mu_K(t)=Q^{-d}\sum_\ell
       \widehat\mu_K(\ell)e^{{\rm i}\theta_\ell^Tt},\qquad
 \mathcal Z(y)=\sum_{t\in[-K,K]^d}\mu_K(t)\,g(y-Ct).
 \label{eq:fourier-inverse}
\end{equation}
These are finite exact sums. Complex terms combine into nonnegative real coefficients in exact arithmetic.
The real balanced amplitudes need not lie on any numerical lattice: only the \emph{integer signatures} are
transformed. The physical map $C$ and Gaussian noise remain unchanged.

Using a smaller grid without an additional argument would merge different admissible sums modulo $Q$.
Their physical Gaussian likelihoods need not agree, so that aliasing would be an approximation.
Likewise, quantizing the real amplitudes before an FFT is unnecessary here and would change the encoder.

\clearpage
\subsection{A count marker makes this a complete PUPE decoder}
Evidence alone is insufficient. For a target message $w_*$ introduce a formal variable $z$ whose power counts
copies of that message, and form
\begin{equation}
 F_\ell(u,z)=\prod_{w\in\calW}
 \left[\sum_{j=0}^K\beta_w(j)u^j
 z^{j\ind\{w=w_*\}}e^{-{\rm i}j\theta_\ell^Tt_w}\right],
 \label{eq:fourier-marker}
\end{equation}
with $\beta_w(j)=1/j!$ for IID sampling and $\beta_w(j)=\ind\{j\leq1\}$ for distinct sampling.
Terms with degree in $u$ greater than $K$ can be discarded exactly.
For IID, the same cardinality coefficient can be written
\[
 [u^K]F_\ell(u,z)=
 \frac1{K!}\left(\sum_w z^{\ind\{w=w_*\}}
                      e^{-{\rm i}\theta_\ell^Tt_w}\right)^K.
\]
The coefficient of $z^j$ records states with exactly $j$ copies of $w_*$. Applying the inverse transform and
the Gaussian weight gives
\begin{align}
 A^{\rm F}_{w_*,j}(y)&=
 \sum_{t\in[-K,K]^d}g(y-Ct)\,
 Q^{-d}\sum_\ell [u^Kz^j]F_\ell(u,z)e^{{\rm i}\theta_\ell^Tt},\\
 q_{w_*,j}&=A^{\rm F}_{w_*,j}/\sum_{a=0}^K A^{\rm F}_{w_*,a}.
 \label{eq:fourier-query}
\end{align}
These are exact-count numerators. Sum the probabilities over $j$ as before to obtain the count tails. Repeating the marked query over full messages completes the Bayes decoder.

\subsection{What changes computationally}
The positive-measure algorithm stores reachable sums. This representation stores frequency samples on a full
grid and uses powers or coefficient products. The transform can remove repeated convolution work when the
effective signature dimension is small. But the direct grid has $(2K+1)^d$ entries.
For the current $P=64$, $|\mathcal G|=29$ representation, $d=1856$: a dense grid is plainly infeasible.
Sparsity of the reachable integer states is often far more valuable than a fast dense transform.

It is still a useful exact alternative because it states precisely what a spectral or low-rank representation
would have to preserve. Reducing the number of retained frequencies, introducing a rank cutoff, or merging
integer states requires an error analysis; none is implied by the Fourier identity itself.
The finite certificate checks a two-dimensional unit-energy alphabet under both priors, including every count
probability. It does not run this grid for the full polar bank.

\subsection{A bound linking Fourier compression to decision certification}
Let $\Gamma_\ell=\sum_t g(y-Ct)e^{{\rm i}\theta_\ell^Tt}$.
Equation \eqref{eq:fourier-inverse} gives $\mathcal Z=Q^{-d}\sum_\ell\widehat\mu_K(\ell)\Gamma_\ell$.
If only frequencies in $\Lambda$ are evaluated, the omitted contribution obeys
\begin{equation}
 |\mathcal Z-\mathcal Z_\Lambda|
 \leq Q^{-d}\sum_{\ell\notin\Lambda}|\widehat\mu_K(\ell)|\,|\Gamma_\ell|.
 \label{eq:fourier-bound}
\end{equation}
For IID draws, write $\psi_\ell=\rho_\ell/M$. Then $|\psi_\ell|\leq1$ and
$|\widehat\mu_K(\ell)|=M^K|\psi_\ell|^K/K!$.
Thus modes strictly below unit magnitude are damped by the power $K$, although the grid and its dimension also
matter. Unit-magnitude modes need not decay. Analogous bounds apply to each marked numerator.
These are absolute bounds; small relative or posterior-tail error additionally needs a positive evidence bound.
Computing useful omitted-mode bounds cheaply remains open. The finite identity supplies a possible connection
to Section~\ref{sec:certificate}, not a guarantee that a short Fourier expansion suffices.

\clearpage
\section{Exact elimination with joint separator states}
\label{sec:elimination-exact}
\subsection{Sum variables in a different order}
The polar/CRC circuit, pattern selectors and resource likelihoods define finite factors of the slot posterior.
An alternative exact algorithm eliminates these variables directly instead of building complete-message sum tables.
Let $x$ collect all finite variables and let $\prod_a f_a(x_{\mathcal S_a})$ be an unnormalized weight whose sum
equals the desired evidence. Deterministic circuit factors ensure that every valid slot assignment has exactly
one valid extension to its internal encoding variables. Include full-message exclusion factors for distinct sampling.

To eliminate a variable $x_i$, collect every factor containing it, multiply them, and sum it out:
\begin{equation}
 \widetilde f(x_{\mathcal N_i})
 =\sum_{x_i}\prod_{a:i\in\mathcal S_a}f_a(x_{\mathcal S_a}),
 \qquad
 \mathcal N_i=\left(\bigcup_{a:i\in\mathcal S_a}\mathcal S_a\right)\setminus\{i\}.
 \label{eq:eliminate}
\end{equation}
Replace those factors by $\widetilde f$ and continue. Distributivity proves that each step preserves the
remaining marginal exactly. A joint table over $\mathcal N_i$ records everything the eliminated variable can
communicate to the rest of the graph. It cannot generally be replaced by separate one-variable tables.

For example, summing the middle variable in a chain creates
\[
 \widetilde f(x_1,x_3)=\sum_{x_2}f(x_1,x_2)h(x_2,x_3),
\]
which usually couples its two neighbors.
Eliminating an endpoint first avoids creating this larger interaction at that step.
The final sum is the same, but the intermediate cost differs.

This is the exact variable-elimination/junction-tree framework \cite{wainwright}.
In tensor language, the factors are arrays and the summed variables are contracted indices.
These descriptions refer to the same sum; calling it a tensor contraction does not by itself make it cheaper.

\subsection{Carry count information through elimination}
For one target message $w_*$, attach a marker $z^{\ind\{w_k=w_*\}}$ to every slot. The final polynomial is
\begin{equation}
 P_{w_*}(z)=
 \sum_W\pi(W)\,g\!\left(y-\sum_k\phi(w_k)\right)
             \prod_{k=1}^K z^{\ind\{w_k=w_*\}}
 =\sum_{j=0}^K A_{w_*,j}z^j.
 \label{eq:elimination-marker}
\end{equation}
Run \eqref{eq:eliminate} with polynomial-valued factors, adding coefficients when summing and convolving them
when multiplying. Degree at most $K$ suffices because there are only $K$ markers.
Then $q_{w_*,j}=A_{w_*,j}/P_{w_*}(1)$, followed by the established tail allocation.
The marker is an exact full-message equality test, not an independent decision about each payload bit.

\clearpage
\subsection{What makes an elimination order tractable}
The main cost is the number of joint assignments in the largest intermediate factor.
For binary variables, retaining $b$ boundary bits means as many as $2^b$ entries; pattern variables have larger
domains. Polynomial markers add a count dimension and polynomial arithmetic, but do not remove this dependence.
The smallest achievable maximum boundary size is related to graph treewidth \cite{wainwright}.
The relevant graph is the chosen factorization, including CRC, pattern consistency and any exclusion constraints.

This gives a precise version of the earlier polar question: find an elimination order whose boundaries preserve
cross-user pairing while remaining small. Ordinary local loopy BP does not retain these full joint boundaries.
The XOR counterexample in Section~\ref{sec:polar-exact} rules out one proposed smaller state, not every
lossless boundary representation. An exact rank factorization of a boundary table could also save storage;
discarding nonzero rank components would introduce an approximation.

\subsection{Conditioning can reveal exact separations}
One exact strategy enumerates a chosen set of variables first, eliminates the rest conditional on them, and
finally sums the conditional answers. If all patterns $(p_1,\ldots,p_K)$ are fixed, connect two slots whenever
their occupied resource sets overlap. Different connected components use disjoint observations.
The Gaussian likelihood then factors by component, apart from an observation-only factor on unused resources.
Within a component, use full joint polar inference. Same-pattern users belong to the same component.

For IID payloads the conditional sum factors immediately. Distinct-message exclusions occur only among slots
with the same pattern and therefore remain inside a component. Summing the \emph{unnormalized} answers over
all pattern assignments restores the full posterior. In particular, distinct sampling must not assign equal
normalized weight to every pattern assignment: they can admit different numbers of valid payload assignments.

If the largest conditional component has only two or three slots, then two- or three-user exact subsolvers can
suffice inside that conditional problem, regardless of total $K$. This is a structural separation, not a claim
that more distant users have negligible influence. Enumerating the patterns may still cost $P^K$.

For the selected $n=256,N_c=32$ banks with $P=64,256$ and seeds 3201--3203, the support-overlap graph on the
entire pattern bank is connected, with about $98.9$--$99.3\%$ of possible edges.
Thus disjoint-support decomposition alone does not split that bank. This read-only structural check is neither
a measured elimination width nor a proof against finer conditional or polar-circuit reductions.

\textbf{Exact computational endpoint.} Eliminate finite variables in any complete order while preserving
joint tables and count markers. The result is the same optimal multiset decoder as before.
The open algebraic task is to find a small exact boundary or a favorable conditioning scheme for this encoder.
The certificate checks the full count law under two elimination orders on the tiny polar fixture, and a
separate chain where a different order changes the largest table from 9 to 27 entries.

\clearpage
\section{An exact expansion in connected groups of users}
\label{sec:connected-exact}
\subsection{Separate isolated evidence from pair interactions}
The Gaussian log-likelihood already contains only individual and pairwise terms, as in \eqref{eq:interactions}.
To expose what those pairs do collectively, define the normalized one-user reference
\begin{equation}
 \nu_y(w)=\frac{e^{y^T\phi(w)/\sigma^2}}{Z_1(y)},\qquad
 Z_1(y)=\sum_u e^{y^T\phi(u)/\sigma^2}.
 \label{eq:one-user-reference}
\end{equation}
For unit-energy messages and a uniform prior, this is the exact one-user posterior, with patterns summed as part
of the message alphabet. It is a reference distribution for the following identity; the true $K$-slot posterior
is not a product of these marginals.

Define the pair factor and its departure from one:
\begin{equation}
 b(u,v)=
 \begin{cases}
 e^{-\phi(u)^T\phi(v)/\sigma^2},&\text{IID},\\
 \ind\{u\ne v\}e^{-\phi(u)^T\phi(v)/\sigma^2},&\text{distinct},
 \end{cases}
 \qquad f(u,v)=b(u,v)-1.
 \label{eq:mayer-factor}
\end{equation}
Then, exactly,
\[
 \Pr(W\mid y)\ \propto\
       \prod_{i=1}^K\nu_y(w_i)\prod_{i<j}b(w_i,w_j).
\]
The distinct model is included by a zero pair factor on duplicate messages.
In either model $b$ is nonnegative, but $f$ can have either sign.

\subsection{Expand the finite product and group connected terms}
For a target message $w_*$ let $J(W)=\sum_i\ind\{w_i=w_*\}$ and define
\begin{equation}
 \mathcal B_k(z)=
 \E_{\nu_y^{\otimes k}}\!\left[z^{J(W)}\prod_{i<j}(1+f(w_i,w_j))\right],
 \qquad \mathcal B_0(z)=1.
 \label{eq:graph-partition}
\end{equation}
Expanding the product chooses a subset of the $\binom k2$ pairs. Draw each chosen pair as an edge between
its slot labels. A term is thus a graph; its connected components partition the slots.
Different components involve disjoint variables, so their sums factor under the product reference.

For $m\geq1$, collect all connected graphs on $m$ labeled slots:
\begin{equation}
 \mathcal C_m(z)=
 \E_{\nu_y^{\otimes m}}\!\left[
 z^{J(W)}
 \sum_{\substack{G\ {\rm connected}\\{\rm on}\ \{1,\ldots,m\}}}
       \prod_{(i,j)\in E(G)}f(w_i,w_j)\right].
 \label{eq:connected-block}
\end{equation}
The singleton graph has no edges, so $\mathcal C_1(z)=\sum_w\nu_y(w)z^{\ind\{w=w_*\}}$.
A connected group means its users are linked by pair interactions. It does not introduce a new many-body
physical channel term.

\clearpage
\subsection{An exact recursion over connected blocks}
In any graph, slot 1 belongs to one connected component of size $m$.
Choose its other $m-1$ slots in $\binom{k-1}{m-1}$ ways, sum the connected component, and sum an arbitrary
graph on the remaining slots. Exchangeability makes each choice have the same weight. Hence
\begin{equation}
 \boxed{\mathcal B_k(z)=
 \sum_{m=1}^k\binom{k-1}{m-1}
                 \mathcal C_m(z)\mathcal B_{k-m}(z).}
 \label{eq:connected-recursion}
\end{equation}
For example,
\[
 \mathcal B_2=\mathcal C_1^2+\mathcal C_2,\qquad
 \mathcal B_3=\mathcal C_1^3+3\mathcal C_1\mathcal C_2+\mathcal C_3.
\]
The connected three-user term includes three two-edge chains and one triangle:
\[
 \mathcal C_3(z)=\E_{\nu_y^{\otimes3}}\!
 \left[z^{J(W)}
 \left(f_{12}f_{13}+f_{12}f_{23}+f_{13}f_{23}+f_{12}f_{13}f_{23}\right)\right].
\]
Equivalently, as a formal series whose coefficients are needed only through degree $K$,
\[
 \sum_{k=0}^K\frac{t^k}{k!}\mathcal B_k(z)
 =\exp\!\left(\sum_{m=1}^K\frac{t^m}{m!}\mathcal C_m(z)\right)\pmod{t^{K+1}}.
\]
Connected groups are the coefficients of the logarithm of this generating function.
All dependence was generated by pair factors; retaining only two-user connected terms is not generally exact.

\subsection{Recover the multiset decision}
The one-user normalizers and the observation-only Gaussian factor do not depend on $J(W)$.
Consequently
\begin{equation}
 \boxed{q_{w_*,j}=
 \frac{[z^j]\mathcal B_K(z)}{\mathcal B_K(1)}.}
 \label{eq:connected-query}
\end{equation}
For IID sampling,
$H_K(y)=e^{-(\|y\|^2+K)/(2\sigma^2)}Z_1(y)^K\mathcal B_K(1)$;
replace $H_K$ by $D_K$ for distinct sampling.
Compute the connected polynomials up to $K$, apply \eqref{eq:connected-recursion}, extract every count law,
and perform the same optimal tail allocation. This is another exact evaluation of the required quantities.

\subsection{Exact expansion and the role of unevaluated terms}
The finite expansion is always an identity; it requires no convergence assumption.
Its naive evaluation is expensive: $\mathcal C_m$ sums message assignments and connected graphs, up to
$M^m2^{\binom m2}$ combinations before exploiting symmetry. Signed terms can also cancel strongly.
The method is related to the classical Mayer expansion; convergence results for other particle systems
require their own hypotheses and are not imported here \cite{pulvirenti}.

Exact inference uses the complete expansion or proves that unevaluated contributions cannot change the
Bayes decision. A partial connected polynomial need not have nonnegative coefficients or preserve all source
exclusions. The next section derives positive bounds on the unevaluated contributions, allowing exact
certification without treating a partial signed sum as a posterior. Assumptions that discard the remainder
without such a certificate belong to Report 9.

\clearpage
\section{Bound the unevaluated connected groups}
\label{sec:connected-bound}
\subsection{Weight interactions by the observation, not just their worst case}
The connected identity suggests a further exact procedure: evaluate small connected blocks, bound all others,
and stop only if the multiset decision is certified. This requires a finite remainder bound; we do not assume
that the first two or three blocks suffice.

Use the same $\nu_y$ and $f$ as above. Define
\begin{equation}
 F(u,v)=|f(u,v)|,\qquad
 \epsilon=\max_{u,v}F(u,v),\qquad
 \alpha=\max_u\sum_v\nu_y(v)F(u,v).
 \label{eq:weighted-interaction}
\end{equation}
The quantity $\alpha\leq\epsilon$ measures the largest average interaction with a message drawn from the
one-user reference. It depends on the observation. It is not a physical noise variance or a fitted parameter.

For each block size $m$, define the positive connected sum
\[
 A_m=\E_{\nu_y^{\otimes m}}
       \sum_{G\ {\rm connected}}\prod_{(i,j)\in E(G)}F(w_i,w_j),\qquad A_1=1.
\]
If $\|P\|_1$ denotes the sum of absolute coefficients of a polynomial, then
$\|\mathcal C_m\|_1\leq A_m$. The marker merely sorts assignments by their count, so taking absolute values
and summing coefficients is bounded by the unmarked positive graph sum.

Every connected graph contains a spanning tree. Summing over all trees may overcount it, which is safe for
a positive bound. For a fixed tree, bound the factors from optional non-tree edges by $1+\epsilon$.
There are $\binom m2-(m-1)=(m-1)(m-2)/2$ such edges.
Average the tree-edge product by removing one leaf at a time: the sum over a leaf's message, conditional
on its parent, is at most $\alpha$. There are $m-1$ leaves removed before the last vertex.
Finally there are $m^{m-2}$ labeled spanning trees, as counted by their length-$(m-2)$ Pr\"ufer sequences.
Consequently
\begin{equation}
 \boxed{A_m\leq a_m:=
 m^{m-2}\alpha^{m-1}(1+\epsilon)^{(m-1)(m-2)/2}\quad(m\geq2),\qquad a_1=1.}
 \label{eq:tree-majorant}
\end{equation}
This is a finite inequality for the declared decoder, including the distinct-message exclusion factor.
It imports no high-temperature or infinite-volume convergence theorem.
Smaller valid positive bounds, or exact values of $A_m$, can replace individual $a_m$.

Computing $\alpha$ directly still involves the alphabet and its pair interactions. A cheap upper bound on it
is safe, but may be uninformative. In the distinct model, $F(u,u)=1$ even at arbitrarily large noise variance;
therefore hard exclusions cannot be treated as disappearing weak interactions.

\clearpage
\subsection{Sum the weighted trees exactly instead of using their worst root}
There is a stronger finite majorant with a computable structure. Define
\[
 Q_m=\E_{\nu_y^{\otimes m}}\sum_{T\ {\rm labeled\ tree}}\prod_{(i,j)\in E(T)}F(w_i,w_j),
 \qquad Q_1=1.
\]
Retain this exact positive tree sum in the spanning-tree argument, instead of replacing each leaf average
by $\alpha$. Then
\begin{equation}
 A_m\leq \overline a_m:=Q_m(1+\epsilon)^{(m-1)(m-2)/2}
 \leq a_m.
 \label{eq:weighted-tree-majorant}
\end{equation}
This alone does not help unless $Q_m$ can be obtained more cheaply than enumerating its trees and messages.

Fix the root's message to $w$. Removing that root separates a labeled tree into an unordered collection of
child subtrees. Each child with message $v$ contributes its edge weight $F(w,v)$ and source weight $\nu_y(v)$.
Choosing the disjoint label sets of the children gives the factorial coefficients of an exponential.
Thus the rooted exponential generating functions obey the formal identity
\begin{equation}
 T_w(z)=z\exp\!\left(\sum_v F(w,v)\nu_y(v)T_v(z)\right),\qquad
 Q_m=(m-1)!\sum_w\nu_y(w)[z^m]T_w(z).
 \label{eq:rooted-tree-series}
\end{equation}
Only coefficients through degree $K$ are needed. This is finite polynomial arithmetic, with no convergence
condition and no large-$K$ limit.

Here is an explicit coefficient algorithm. Write $t_{w,m}=[z^m]T_w$,
$\ell_{w,j}=\sum_vF(w,v)\nu_y(v)t_{v,j}$, and
$e_{w,n}=[z^n]\exp(\sum_{j\geq1}\ell_{w,j}z^j)$.
Initialize $e_{w,0}=t_{w,1}=1$. Differentiating the exponential gives
\[
 e_{w,n}=\frac1n\sum_{j=1}^n j\ell_{w,j}e_{w,n-j},\qquad t_{w,n+1}=e_{w,n}.
\]
At degree $n$, all required $t_{v,n}$ are already known. Compute the new vector
$\ell_{\cdot,n}=F(\nu_y\,t_{\cdot,n})$, then the next coefficients.
Through order $K$ this needs $K-1$ matrix-vector products, $O(K^2M)$ additional scalar operations,
and $O(KM)$ coefficient storage. The exact polar operator in Section~\ref{sec:pair-contraction} can supply
those products.

For a check, put $d=F\nu_y$. The first cases are
\[
 Q_2=\nu_y^Td,\qquad Q_3=3\sum_w\nu_y(w)d_w^2.
\]
The old three-user bound used $3\max_wd_w^2$ instead. Constant interaction $F=c$ recovers
$Q_m=m^{m-2}c^{m-1}$. Tests also enumerate all labeled trees through five vertices for a nonconstant kernel.
Substitute $\overline a_m$ for $a_m$ in the following remainder recurrence. It is never weaker in exact
arithmetic, but optional-edge overcounting, the $1+\epsilon$ factors and signed cancellation can still make
the resulting bound loose.

\clearpage
\subsection{Turn block bounds into a bound on the complete count polynomial}
Let $\widetilde{\mathcal B}^{(b)}_K$ be the recurrence in \eqref{eq:connected-recursion} evaluated with the
exact $\mathcal C_m$ through $m=b$ and zero blocks above $b$. Define positive scalar recurrences
\[
 R_0^{(b)}=1,\qquad
 R_k^{(b)}=\sum_{m=1}^{\min(k,b)}\binom{k-1}{m-1}a_mR_{k-m}^{(b)}.
\]
Then
\begin{equation}
 \boxed{\|\mathcal B_K-\widetilde{\mathcal B}^{(b)}_K\|_1
 \leq\delta_b:=R_K^{(K)}-R_K^{(b)}.}
 \label{eq:connected-remainder}
\end{equation}
To see why, expand the exact answer over partitions into connected components. The difference contains exactly
the partitions with at least one component larger than $b$. The polynomial norm is submultiplicative, so
each such partition is bounded by the product of its $a_m$ values. The positive recurrence sums those products,
and the subtraction removes the partitions already retained. It gives $\delta_K=0$.

An independent coarse bound is
$\sum_{\ell=b}^{\binom K2}\binom{\binom K2}{\ell}\epsilon^\ell$:
an omitted graph must contain at least $b$ edges. Take the smaller of any valid bounds.
The weighted tree bound can improve this crude estimate, but neither is promised small for the present encoder.
Both can grow rapidly with $K$ or with strong interactions.

\subsection{Certify a decision even when an intermediate polynomial has negative coefficients}
For a queried message write
$\widetilde N_{w,j}=\sum_{\ell\geq j}[z^\ell]\widetilde{\mathcal B}^{(b)}_K(z)$ and
$\widetilde Z=\widetilde{\mathcal B}^{(b)}_K(1)$. Then
\[
 |N_{w,j}-\widetilde N_{w,j}|\leq\delta_b,\qquad
 |\mathcal B_K(1)-\widetilde Z|\leq\delta_b.
\]
All queries use the same positive normalizer $\mathcal B_K(1)$.
Thus valid unnormalized intervals are
\[
 L_{w,j}=\max(0,\widetilde N_{w,j}-\delta_b),\qquad
 U_{w,j}=\max(0,\widetilde N_{w,j}+\delta_b),
\]
with $Z_-=\max(0,\widetilde Z-\delta_b)$ and $Z_+=\widetilde Z+\delta_b>0$.
One may tighten $U_{w,j}$ by $Z_+$. Normalized tails obey
\[
 \frac{L_{w,j}}{Z_+}\leq T_{w,j}\leq
 \begin{cases}\min(1,U_{w,j}/Z_-),&Z_->0,\\1,&Z_-=0.\end{cases}
\]
The selected-versus-unselected interval test of Section~\ref{sec:certificate} applies directly, before
normalization. If it separates the count tails, the resulting action is exactly Bayes optimal even though not
all connected blocks were evaluated. If it does not, increase the block size, improve the bounds, or use another
exact representation. Full evaluation remains the finite endpoint.

Negative coefficients in $\widetilde{\mathcal B}$ are not probabilities and must not be normalized as if they
were. They are nevertheless permissible as a partial signed calculation accompanied by a rigorous error interval.
This distinction connects the earlier truncation counterexample to an exact certification procedure.
Floating point evaluation of the bounds still needs controlled rounding for machine-certified guarantees.

\textbf{A finite early-stop example.} Take the four unit waveforms $\{e_1,e_2,-e_1,-e_2\}$, IID $K=4$,
$y=(0.7,-0.2)$ and $\sigma^2=5$. Propose one copy of each message. The computed bounds are:
\begin{center}
\begin{tabular}{@{}rrrr@{}}\toprule
Block size $b$ & Remainder $\delta_b$ & Selected lower & Next-copy upper\\\midrule
1 & 0.832126 & 0 & 1.155913\\
2 & 0.185467 & 0.527514 & 0.499847\\
3 & 0.031415 & 0.671503 & 0.344105\\\bottomrule
\end{tabular}
\end{center}
At $b=2$ the selected lower bounds already exceed every next-copy upper bound.
The two-user calculation, together with the unevaluated-group bound, therefore certifies the exact four-user
Bayes action. Exhaustive inference independently agrees. This is a constructed finite illustration, not a
performance result for the polar bank or a claim that size two always suffices.

\clearpage
\section{Compute only the comparisons needed by the decision}
\label{sec:decision-computation}
\subsection{Compare count tails jointly instead of bounding each one separately}
For a proposed allocation $\widehat a$, consider removing its last copy of message $u$ and adding the next copy
of message $v\ne u$. Put $j=\widehat a_u>0$ and $\ell=\widehat a_v+1$. The unnormalized expected-match advantage
of retaining that selected copy is
\begin{equation}
 D_{u,j;v,\ell}
 =N_{u,j}-N_{v,\ell}
 =\sum_a\omega(a)\big[\ind\{a_u\geq j\}-\ind\{a_v\geq\ell\}\big].
 \label{eq:decision-contrast}
\end{equation}
All such advantages being nonnegative is equivalent to the allocation being Bayes optimal.
Necessity follows because a negative one gives an improving one-copy exchange.
For sufficiency, the diminishing count tails ensure that these last-selected versus next-unselected comparisons
also order every selected occurrence above every unselected occurrence.
Comparisons within the same message are already nonnegative by tail monotonicity.

The earlier interval certificate bounds each tail separately. It can lose cancellation: an unexplored state
in which both indicators equal one contributes to both tails but contributes zero to their difference.
For a count-tree frontier node $b$, let $E_b^-$ mean that a completion with
$a_u<j$ and $a_v\geq\ell$ is possible. With the same upper evidence mass $U_b$ as before, a sharper bound is
\begin{equation}
 \boxed{D_{u,j;v,\ell}\ \geq\
 \sum_{a\in\mathcal E}\omega(a)
       \big[\ind\{a_u\geq j\}-\ind\{a_v\geq\ell\}\big]
 -\sum_{b\in\mathcal F:E_b^-\ {\rm possible}}U_b.}
 \label{eq:direct-contrast-bound}
\end{equation}
If the negative event is impossible, every completion has nonnegative contrast and needs no penalty.
If it is possible, its total negative contribution is at worst $-U_b$.
Positive lower bounds on node mass could improve this further, but are not required here.

Feasibility is finite arithmetic on count ranges. Fix the prefix counts, impose
$a_u\leq j-1$, $a_v\geq\ell$, nonnegativity, total $K$, and the distinct-source upper caps when applicable.
The remaining intervals admit a completion exactly when their lower endpoints sum to at most $K$ and their
upper endpoints sum to at least $K$, with no empty interval. No transmitted-message oracle is used.

At the same frontier, \eqref{eq:direct-contrast-bound} is at least as strong as subtracting independent tail
intervals. It penalizes a subset of the nodes that could increase the unselected tail.
Certify the action once every required contrast lower bound is nonnegative; otherwise continue exact refinement.
This preserves common uncertainty rather than estimating the entire posterior more accurately.

The same reasoning helps the connected expansion. Its positive omitted-graph bound $\delta_b$ controls the
integral of any observable bounded in magnitude by one. A difference of two indicators has this property.
Hence its contrast error is at most $\delta_b$, rather than the $2\delta_b$ obtained by separately bounding
two tails. This does not guarantee a useful bound, but avoids one unnecessary loss.

\clearpage
\subsection{A controlled finite audit of bound strength}
Use the existing balanced polar fixture $B=4$, $P=4$, $V=4$, $N_c=8$, CRC4, $n=12$, $K=3$.
For each source law and each $\sigma^2\in\{0.05,0.3,1\}$, generate four observations with random seeds
9100--9103: first draw the three messages under that source law, then draw Gaussian noise.
There are 24 observations in total, with 816 IID or 560 distinct count hypotheses per observation.
This is a small computational audit, not an error-rate study of the research-scale encoder.

Within each observation all methods use the same count-tree order, obtained by decreasing matched score
$y^T\phi(w)$ with a deterministic tie break. Refine the frontier node with largest upper evidence mass.
Propose the allocation using the evaluated lower tails. Compare three bounds:
\begin{enumerate}
\item \textbf{Ball:} the original norm bound.
\item \textbf{Box:} also bound each completion coordinate by its smallest and largest possible value.
      For IID these are $h\min_w\phi_r(w)$ and $h\max_w\phi_r(w)$; for distinct sampling, sum the $h$ smallest
      or largest coordinate values of different remaining messages. Distance to this containing box is a
      valid squared-distance lower bound. Keep the stronger of it and the ball bound.
\item \textbf{Dual:} additionally take up to eight conditional-gradient steps in the completion hull, starting
      at the feasible point $h$ times the remaining alphabet mean. Each step uses the exact remaining-alphabet
      support and evaluates \eqref{eq:completion-dual-bound}. Keep the largest valid bound seen.
      The convex solve need not converge.
\end{enumerate}
The resulting mean percentages of full count hypotheses evaluated before certification are:
\begin{center}
\begin{tabular}{@{}llrrrr@{}}\toprule
Source & $\sigma^2$ & Ball / intervals & Box / intervals & Dual / intervals & Dual / contrasts\\\midrule
IID & 0.05 & 59.68 & 23.53 & 3.25 & 3.22\\
IID & 0.30 & 87.44 & 55.51 & 25.06 & 23.62\\
IID & 1.00 & 91.54 & 67.25 & 43.69 & 43.11\\
Distinct & 0.05 & 63.13 & 20.31 & 1.34 & 0.98\\
Distinct & 0.30 & 94.60 & 57.19 & 24.29 & 22.46\\
Distinct & 1.00 & 95.80 & 71.38 & 52.72 & 51.03\\\bottomrule
\end{tabular}
\end{center}
All 144 combinations of observation, bound and certificate agree with exhaustive Bayes inference to the stated
floating tolerances. Both certificate variants were tested with all three bounds; the table shows the most
informative comparisons. The contrast improvement is modest here; the main reduction comes from stronger
completion geometry. We do not extrapolate these four-observation averages to other loads or message lengths.

\clearpage
\subsection{Count the work needed to obtain the savings}
The dual bounds are not free. For the dual/contrast method, the mean work counts in the same audit are:
\begin{center}
\begin{tabular}{@{}llrrr@{}}\toprule
Source & $\sigma^2$ & Evaluated hypotheses & Support calls & Individual message scores\\\midrule
IID & 0.05 & 26.25 & 322.00 & 3,682.00\\
IID & 0.30 & 192.75 & 1,888.00 & 15,565.75\\
IID & 1.00 & 351.75 & 3,288.25 & 22,894.00\\
Distinct & 0.05 & 5.50 & 192.00 & 2,318.00\\
Distinct & 0.30 & 125.75 & 1,720.50 & 13,627.00\\
Distinct & 1.00 & 285.75 & 3,194.75 & 20,546.25\\\bottomrule
\end{tabular}
\end{center}
A support call scores every message in the node's remaining alphabet; the last column counts those scores.
These are operation counts, not comparable units of elapsed time. Node construction, likelihood evaluation
and certificate comparisons also cost work. In particular, this audit does not show that the reference search
is faster than vectorized exhaustive inference on such tiny systems.

It does establish a more specific fact: the final Bayes decision can often be certified without evaluating most
count hypotheses, using the actual polar fixture rather than only a constructed four-message example.
Whether this pays computationally depends on cheaper support bounds, reuse between nodes, and the decision
margin of the observation. The bound and the exact support interface now give explicit objects to investigate.

\subsection{The current connected majorant does not yet pass the same check}
On these same 24 observations, compute the exact one- and two-user connected blocks and bound the omitted
three-user block by $3\alpha^2(1+\epsilon)$. Use even the stronger direct-contrast error bound just described.
It certifies none of the 24 outputs. Relative to the independently computed exact normalizer
$\mathcal B_3(1)$, the remainder-bound ranges are:
\begin{center}
\begin{tabular}{@{}llrr@{}}\toprule
Source & $\sigma^2$ & Minimum $\delta_2/\mathcal B_3(1)$ & Maximum\\\midrule
IID & 0.05 & $1.88\,10^{15}$ & $7.25\,10^{31}$\\
IID & 0.30 & $228.36$ & $13{,}240.39$\\
IID & 1.00 & $0.677$ & $1.610$\\
Distinct & 0.05 & $1.93\,10^8$ & $6.87\,10^{24}$\\
Distinct & 0.30 & $76.73$ & $8{,}959.65$\\
Distinct & 1.00 & $1.137$ & $3.325$\\\bottomrule
\end{tabular}
\end{center}
This is a failure of this sufficient certificate on these cases. It neither proves a large actual truncation
error nor excludes sharper connected bounds or a different exact decomposition.
The earlier four-message early-stop example remains valid, but is not evidence of useful certification here.

Reproduce the complete audit with \texttt{uv run python -m tests.polar\_exact\_computation}.
Generated audit data are kept outside the repository; the derivation, procedure, seeds and reported tables
are retained here. All numerical certificates are checked against exhaustive reference decisions; formal
machine certification would additionally require controlled-rounding bounds.

\clearpage
\section{Reuse valid bounds and charge only the adverse prior mass}
\label{sec:bound-reuse}
\subsection{A support direction can serve many different subtrees}
The finite audit exposed repeated support scoring as a substantial cost. The dual inequality does not require
a direction tailored to the current residual. Any previously computed direction remains valid.

For a count-tree node with prefix depth $d$, residual $r$ and remaining count $h$, let
$S_d$ be the remaining message suffix. For each cached direction $u_j$, store its full-alphabet scores
$s_{j,w}=u_j^T\phi(w)$. Precompute
\[
 H_{j,d,h}=
 \begin{cases}
 h\max_{w\in S_d}s_{j,w},&\text{IID},\\
 \text{sum of the $h$ largest scores in $S_d$},&\text{distinct}.
 \end{cases}
\]
At every descendant, the same valid bound is
\begin{equation}
 L_{d,h}(r)=\max\{L_{\rm box},\ \max_j(2u_j^Tr-\|u_j\|^2-2H_{j,d,h})\}.
 \label{eq:cached-support-bound}
\end{equation}
Here $L_{\rm box}$ includes zero and the ball bound as before. No new message score is needed for these
dual trials. Along an edge assigning $c$ copies of word $w$, one may update
$u_j^Tr'=u_j^Tr-cs_{j,w}$ instead of taking a fresh residual-direction dot product.
The small reference computes the dots directly and reports their count separately.

The cached vectors may be poor directions after the residual changes substantially. This weakens the bound
and may increase exact search, but it does not approximate the posterior or invalidate a certified decision.
New directions can be added adaptively if their computational cost is worthwhile.
Suffix maxima require $O(JM)$ preparation; maintaining the top $K$ scores per suffix can use
$O(JMK)$ preparation and storage. The simple finite code sorts suffixes directly; the better preparation bound
describes a straightforward maintained-list implementation, not a timing measurement.

For the controlled comparison below, take up to eight directions visited by the root's conditional-gradient
calculation and reuse only those. Each is scored once on the whole $M=16$ alphabet: all tested roots use eight
directions, giving 128 individual message scores in total. The box bounds, cached trial evaluations and
certificate comparisons still require work. No support directions are learned.

\clearpage
\subsection{Use the prior mass of the event that could reverse a decision}
The previous direct-contrast bound charged the full node mass whenever a negative event was possible.
The model also tells us exactly how much prior mass is assigned to that event.

At node $b$, let $\beta_b(a)$ be the unnormalized prior weights of its completions, including its fixed prefix,
and let $B_b=\sum_a\beta_b(a)$. Suppose the completion bound gives a common likelihood envelope
$g(y-\Phi a)\leq G_b$, so $U_b=B_bG_b$.
For $E^-_{u,j;v,\ell}=\{a_u<j,\ a_v\geq\ell\}$, define
\[
 \rho_b^-=
 \frac{\sum_a\beta_b(a)\ind\{E^-_{u,j;v,\ell}\}}{B_b}.
\]
Then the node's negative contribution has magnitude at most $U_b\rho_b^-$, and
\begin{equation}
 \boxed{D_{u,j;v,\ell}\geq D_{u,j;v,\ell}^{\rm evaluated}
                  -\sum_{b\in\mathcal F}U_b\rho_b^-.}
 \label{eq:prior-weighted-contrast}
\end{equation}
The proof is simply to apply the common envelope only to the adverse subset before summing prior weights.
Since $0\leq\rho_b^-\leq1$, this improves the feasibility-only bound at the same frontier.
It uses a conditional \emph{prior} probability, not a posterior probability or an independence approximation.

For two different, still-unassigned words among $m$ remaining messages, and $h$ remaining transmissions,
the IID conditional count law is
\[
 \Pr_{\rm prior}(a_u=a,a_v=b\mid\text{prefix})
 =\frac{h!}{a!b!(h-a-b)!}\frac{(m-2)^{h-a-b}}{m^h}.
\]
Use $a,b\geq0$, $a+b\leq h$, and $0^0=1$. For distinct sampling it is
\[
 \Pr_{\rm prior}(a_u=a,a_v=b\mid\text{prefix})
 =\frac{\binom{m-2}{h-a-b}}{\binom mh},\qquad a,b\in\{0,1\},
\]
with infeasible binomial coefficients zero. Sum over $a<j,b\geq\ell$.
Already assigned words have deterministic counts; when only one unassigned word is involved, use its
binomial or hypergeometric marginal. A single remaining word has deterministic count $h$.
All required probabilities can be cached by $(m,h)$ and the thresholds, separately for each source law.

One can go further by deriving a likelihood envelope specifically for $E^-$, incorporating its count bounds
into the support calculation. The present result does not need that extra work: it already makes the prior
and the decision criterion improve the existing geometric bound together.

\clearpage
\subsection{What the new controlled comparison establishes}
Use exactly the preceding 24 observations, message orders, refinement rule and exhaustive references.
The table gives mean percentages of count hypotheses evaluated before certification:
\begin{center}
\begin{tabular}{@{}llrrrr@{}}\toprule
Source & $\sigma^2$ & Cached / feasible & Cached / prior & Fresh / feasible & Fresh / prior\\\midrule
IID & 0.05 & 13.54 & 12.90 & 3.22 & 3.00\\
IID & 0.30 & 41.82 & 40.01 & 23.62 & 22.37\\
IID & 1.00 & 54.75 & 54.57 & 43.11 & 43.11\\
Distinct & 0.05 & 14.20 & 11.47 & 0.98 & 0.54\\
Distinct & 0.30 & 45.80 & 42.77 & 22.46 & 20.67\\
Distinct & 1.00 & 62.32 & 61.47 & 51.03 & 51.03\\\bottomrule
\end{tabular}
\end{center}
``Feasible'' uses the direct negative-event feasibility penalty; ``prior'' uses
\eqref{eq:prior-weighted-contrast}. ``Fresh'' takes up to eight new dual steps at each nonterminal node.
The cached/prior variant uses 128 message scores per observation, versus group means of
1,974--22,894 for fresh/prior. It additionally evaluates 930--4,052 cached residual-direction dot products
on average across these groups. Thus fewer support scores are purchased with more branching and bound queries.
All 120 runs, including a box-bound control, attain the exhaustive Bayes gain to the stated floating tolerance.
These operation counts identify a tradeoff; a runtime advantage has not been established.

The weighted-tree recurrence was checked on the same observations. For the omitted three-user block,
the new bound divided by the old bound lies in these ranges:
\begin{center}
\begin{tabular}{@{}llrr@{}}\toprule
Source & $\sigma^2$ & Minimum ratio & Maximum ratio\\\midrule
IID & 0.05 & $1.97\,10^{-7}$ & $1.06\,10^{-4}$\\
IID & 0.30 & 0.0524 & 0.5620\\
IID & 1.00 & 0.5758 & 0.7042\\
Distinct & 0.05 & $1.13\,10^{-6}$ & 0.00920\\
Distinct & 0.30 & 0.2007 & 0.4336\\
Distinct & 1.00 & 0.5365 & 0.6791\\\bottomrule
\end{tabular}
\end{center}
Despite these improvements, the new sufficient bound still certifies \textbf{0 of 24} decisions.
Summing the trees more accurately has not removed the remaining overcounting and worst-edge penalties.
This result gives no justification for treating omitted connected groups as negligible.

Reproduce this audit and the six-bank pair-storage census:
\begin{quote}\small\ttfamily
uv run python -m tests.polar\_exact\_reuse
\end{quote}
The full-alphabet pair contractions, two-user posterior, prior-event probabilities, reused support bounds and
rooted-tree recurrence have separate finite checks. The large-bank table is a structural census, not a measured
large-bank decoder run. Formal machine certificates still require rounding control.

\clearpage
\section{Identifiability on one pattern, through the first exact aliases}
\label{sec:identifiability}
Same-pattern collisions are common enough that they cannot be excluded as an implementation convenience.
For IID messages a tagged user's pattern collides with probability $1-(1-1/P)^{K-1}$.
For uniformly distinct full messages the exact probability is
\[
 1-\frac{\binom{M-V}{K-1}}{\binom{M-1}{K-1}},
\]
with an impossible numerator event assigned probability zero. Different payloads can share the same pattern under
either law.

\subsection{Separate ambiguity of the encoder from difficulty of decoding}
Suppose exactly two users share a \emph{known} pattern, other interference has been removed, and noise is zero.
Assume every $e_{p,i}>0$. The receiver can invert insertion/scaling:
\begin{equation}
 t=D_p^{-1}S_p^T\big(\phi(p,v)+\phi(p,v')\big)=b(v)+b(v')\in\{-2,0,2\}^{N_c}.
\end{equation}
Let $d=v\oplus v'$. If $\ell_i^T$ denotes row $i$ of $L$, then
\begin{equation}
 t_i=0\ \Longleftrightarrow\ \ell_i^Td=1,\qquad
 t_i/2=(-1)^{\ell_i^Tv}\quad\hbox{when }\ell_i^Td=0.
 \label{eq:pair}
\end{equation}
The zero/nonzero pattern gives all bits of $Ld$. If $\operatorname{rank}_{\GF(2)}L=B_c$, it uniquely determines $d$.
For $d\ne0$, let $L_0(d)$ contain the rows with $\ell_i^Td=0$. The remaining signs give the linear equations
$L_0(d)v=c_0$. Since $L_0(d)d=0$, their rank is at most $B_c-1$.

\begin{equation}
 \boxed{\begin{gathered}
 \text{The unordered pair is uniquely determined for difference }d\ne0\\
 \Longleftrightarrow\quad \operatorname{rank}_{\GF(2)}L_0(d)=B_c-1.
 \end{gathered}}
 \label{eq:rank}
\end{equation}
If the rank equals $B_c-1$, the solutions are exactly $v$ and $v+d$, the two orderings of the same pair.
If the rank is smaller, the solution affine space has more than two points and contains multiple unordered pairs
with the same sum. For $d=0$, $t=2b(v)$ identifies the repeated payload by full rank of $L$.

\textbf{New local certificate.} The current $32\times6$ effective generator with CRC8 has rank six.
For all 63 nonzero $d$, the restricted rank is five. Independently enumerating all
$\binom{64+1}{2}=2{,}080$ unordered payload pairs gives 2,080 distinct sign sums.
Thus two messages on a fixed pattern have no exact sum alias for this selected code. Positive energy shaping
does not change this conclusion: it is an invertible diagonal map on the pattern support.

\textbf{What it does not establish.} Small energies can amplify noise after inversion. In a noisy mixture,
recognizing exact zero coordinates is not a decoder. Unknown patterns, cross-pattern alternatives and distance
to competing sums require additional analysis. The following extension settles the first same-pattern aliases,
but does not justify a one-message-per-pattern model or certify the full $K$-user encoder.

More generally, if two allowed count vectors satisfy $\Phi a=\Phi a'$, the likelihood cannot distinguish them at
any SNR. Near aliases determine difficult noisy comparisons. Before assigning every error to approximate inference,
use exact tiny-system posteriors and sum-distance diagnostics to distinguish these cases.

\clearpage
\subsection{Three messages remain identifiable; four do not}
Continue with the selected $B_c=6,N_c=32$, CRC8 code and one known pattern.
An exhaustive \emph{integer} sign-sum calculation gives:
\begin{center}
\begin{tabular}{@{}rrrr@{}}\toprule
Payload count & Unordered multisets & Distinct sign sums & Ambiguous sum classes\\\midrule
2 & 2,080 & 2,080 & 0\\
3 & 45,760 & 45,760 & 0\\
4 & 766,480 & 766,052 & 404\\\bottomrule
\end{tabular}
\end{center}
At four users, 392 classes contain two multisets and 12 contain four; all other classes are singletons.
Every one of the 832 ambiguous multisets contains four distinct payloads. The first ambiguity therefore occurs
at four users under \emph{both} source laws, not merely because IID sampling permits repeated messages.

One explicit equality, using decimal indices for the repository's six-bit payload vectors, is
\begin{equation}
 b(12)+b(15)+b(57)+b(58)
 =b(20)+b(23)+b(33)+b(34).
 \label{eq:four-user-alias}
\end{equation}
It is preserved by every fixed $A_p$, including all positive balanced energy choices in this investigation.
It is not a floating point near-collision.

\subsection{Why the equality holds: characters cannot distinguish these two cosets}
Let $U$ be a binary linear subspace. Orthogonality on that subspace gives
\[
 \sum_{v\in a+U}\chi_g(v)
 =\chi_g(a)\sum_{u\in U}\chi_g(u)
 =\begin{cases}|U|\chi_g(a),&g\in U^\perp,\\0,&g\notin U^\perp.\end{cases}
\]
For the zero case, pair each $u$ with $u+u_0$ where $g^Tu_0=1$; their signs cancel.
Consequently the two cosets $a+U$ and $a+d+U$ have identical observed character sums if and only if
\begin{equation}
 g^Td=0\quad\text{for every }g\in\mathcal G\cap U^\perp.
 \label{eq:coset-alias}
\end{equation}
Here $U=\{0,3,53,54\}$ is the span of the binary vectors 3 and 53, $a=12$ and $d=24$.
The two cosets are exactly the sets in \eqref{eq:four-user-alias}.
Direct binary checks of the generator verify \eqref{eq:coset-alias}.
This is an algebraic explanation of the alias, not just an enumeration result.

More generally let $H_{\mathcal G}$ have entries $(H_{\mathcal G})_{g,v}=\chi_g(v)$.
An integer difference $z=a-a'$ is an ambiguity precisely when
$H_{\mathcal G}z=0$ and $\one^Tz=0$, with both count vectors allowed by the source law.
Its positive and negative parts have the same mass. Minimizing that mass over nonzero feasible differences
finds the first possible alias; for this code the minimum is four. Under distinct sampling the entries of a
reduced difference are restricted to $\{-1,0,1\}$.
This connects incomplete Walsh observations to a concrete integer-kernel question.

\clearpage
\subsection{Compute the irreducible risk, with the conditioning stated}
Consider a separate noiseless experiment with exactly four users, all on one \emph{known} pattern, and payloads
drawn uniformly from the 64 possibilities under either source law. Let
$\mathcal A_s=\{a:H_{\mathcal G}a=s,\ \one^Ta=4\}$ be the compatible count class.
An exact observation identifies the class, but its posterior within that class is only
\begin{equation}
 \Pr(a\mid s)=\frac{\beta(a)}{\sum_{a'\in\mathcal A_s}\beta(a')},
 \qquad a\in\mathcal A_s.
 \label{eq:noiseless-class}
\end{equation}
This remains the prior ratio between equal-signal hypotheses at every positive noise variance and every $y$.
More computation cannot distinguish their identical likelihoods.

All members of the ambiguous four-user classes have four distinct payloads, hence equal $\beta=1$ under both
source laws. Within each class their payload sets are pairwise disjoint. A class of size $c$ therefore has
posterior inclusion probability $1/c$ for each represented payload. The best four-occurrence output gains
$4/c$ matches in expectation, giving conditional multiset risk $1-1/c$.
Choosing one entire class member or mixing four supported payloads both attain this Bayes risk.

The total risk numerator over unordered count states is
\[
 392\cdot2\cdot\tfrac12+12\cdot4\cdot\tfrac34=428.
\]
Under IID draws each such four-distinct multiset has probability $4!/64^4$; under distinct draws it has
probability $1/\binom{64}{4}$. Thus the exact Bayes risks in this conditional noiseless experiment are
\begin{equation}
 R_{\rm IID}^{(4,\ {\rm known}\ p)}
 =\frac{4!\,428}{64^4}=0.0006122589\ldots,\qquad
 R_{\rm dist}^{(4,\ {\rm known}\ p)}
 =\frac{428}{\binom{64}{4}}=0.0006736169\ldots.
 \label{eq:four-user-risk}
\end{equation}
These are not the full bank's measured PUPE, nor an estimate of the practical importance of four-way pattern
collisions. The experiment conditions on all four users sharing a known pattern and supplies an exact signal.
Unknown patterns and observation noise introduce additional uncertainty.

The equality nevertheless embeds in larger systems: append the same other messages to both sides of
\eqref{eq:four-user-alias}. Under distinct sampling those extra messages must respect the no-duplicate constraint.
This proves the existence of larger ambiguous frames without claiming that they are common or dominate error.

\textbf{Consequence for decoder research.} Exact inference must preserve the probability mass and count tails of
all compatible explanations. A correct receiver can return a Bayes action while remaining uncertain about the
transmitted multiset. The finding does not reopen encoder optimization here; it records a limitation of the
frozen construction that no proposed exact or approximate decoder should be expected to remove.

\clearpage
\section{Exact reference changes, invariants and decision contrasts}
\label{sec:synthesis}
\subsection{The same dependence appears in two different representations}
Several exact conclusions fit together: continuation-row proportionality, removable unary factors,
message ambiguity and marked decision contrasts share common algebraic structure.
The companion \path{docs/POLAR_DECODING_RESEARCH_LOG.md} records the synthesis and computational limits.
The deductions here concern this model; they are not claims of new general mathematical techniques.

Start with the positive Gaussian pair factor $b(u,v)=e^{-\phi(u)^T\phi(v)/\sigma^2}$. Dividing it by any
positive $a(u)d(v)$ and multiplying the endpoint unary weights by $a$ and $d$ changes no configuration weight.
Thus some apparent interaction can be reassigned exactly to individual-message evidence. The invariant is
\begin{equation}
 \log\frac{b(u,v)b(u',v')}{b(u,v')b(u',v)}
 =-\frac{[\phi(u)-\phi(u')]^T[\phi(v)-\phi(v')]}{\sigma^2}.
 \label{eq:cross-ratio}
\end{equation}
On a full positive rectangular table, every such ratio being one is equivalent to factorization into unary
functions: fix one row and column and reconstruct every entry. Distinct-source zeros cannot be handled by
taking their logarithms; those hard constraints remain unchanged.

Now replace the endpoints by a partial signal $s$ and a legal future signal $t$. The same four-point expression
is $-(s-s')^T(t-t')/\sigma^2$. Its vanishing for every future pair is precisely the continuation-row
proportionality criterion of Section~\ref{sec:exact-reductions}. Thus pair dependence and retained continuation
information are two views of the same interaction. A large Gram entry may include a removable common component;
its magnitude alone does not measure irreducible dependence.

\subsection{Choose the exact reference before judging a connected remainder}
For any vector $m$, put $x_w=\phi(w)-m$ and define
\begin{equation}
 \nu_w^{(m)}\propto e^{[y-(K-1)m]^T\phi(w)/\sigma^2},\qquad
 b^{(m)}(u,v)=e^{-x_u^Tx_v/\sigma^2}.
 \label{eq:centred-reference}
\end{equation}
Keep the original distinct-message exclusions if applicable. Expanding the product's exponent shows that its
unnormalized weight differs from the original by only $e^{-\binom K2\|m\|^2/\sigma^2}$. The posterior is identical.
The physical codewords have not changed; the shifted $x_w$ are an algebraic representation, not a new encoder.

The connected expansion may now use this reference. Its $f^{(m)}=b^{(m)}-1$, weighted tree sums and sufficient
remainder bounds can differ. Under distinct sampling the hard zeros must be included when forming $f$, whose
diagonal remains $-1$. For a fixed target $K$, freeze the chosen reference throughout the block recurrence;
recomputing it for each smaller block size would invalidate the coefficient identity.

A constructed example illustrates the distinction. Four unit vectors share a component of norm
$\sqrt{1-.05^2}$ and have transverse directions $(.05,0),(0,.05),(-.05,0),(0,-.05)$. At $\sigma^2=.3$,
the maximum original $|b-1|$ exceeds $.96$. Centering out their common component reduces it below $.009$
without changing any posterior probability. The actual polar bank has full global variation span; this example
does not establish a similarly useful center there. It shows why a failed sufficient bound at one reference
does not rule out all references.

The positive pair contraction remains available. With $d_w=e^{m^T\phi(w)/\sigma^2}$,
\begin{equation}
 B^{(m)}x=e^{-\|m\|^2/\sigma^2}D_dB(D_dx),\qquad D_d=\diag(d).
 \label{eq:centred-operator}
\end{equation}
This adds only pointwise operations to a $B$ application, including the same distinct-source zeros. But
$|B^{(m)}-1|$ is \emph{not} this diagonal scaling of $|B-1|$. Its old parity tables may need extra state or a
different bound. Improved remainder size must therefore be assessed alongside contraction cost and numerical
stability. Recentring leaves the four-point ratios in \eqref{eq:cross-ratio} unchanged.

\clearpage
\subsection{Known cardinality gives an exact count-factor reparameterization}
Write $A=\Phi^T\Phi/\sigma^2$ and $l=\Phi^Ty/\sigma^2$. For any $\lambda\in\R^M$, put
\[
 A'=A-\lambda\one^T-\one\lambda^T,\qquad l'=l-K\lambda.
\]
For every legal count vector with $\one^Ta=K$,
\[
 l'^Ta-\tfrac12a^TA'a
 =l^Ta-K\lambda^Ta-\tfrac12a^TAa+(\one^Ta)\lambda^Ta
 =l^Ta-\tfrac12a^TAa.
\]
The source prior and posterior are unchanged. $A'$ need not be a physical positive-semidefinite Gram matrix;
it represents the same finite count weights on the known-cardinality domain. This identity can expose
removable pair contributions, but supplies neither a cheap choice of $\lambda$ nor permission to omit the
remaining interactions. The notation does not require constructing a full $M$-by-$M$ table.

\subsection{Joint interfaces control exact regional contraction}
The K2/K3 solvers eliminate internal variables while retaining their joint boundary dependence. The census
in Section~\ref{sec:three-user} shows that two neighbours can expose a whole payload. A smaller region need
not have a smaller interface, and unknown patterns and occupancies remain summation variables. Exact
conditional calculations with true patterns supplied would instead be oracle calculations. Discarding a
boundary dependence is an additional assumption, treated in Report 9.

\clearpage
\subsection{Exact conditional identities inside ambiguous signal classes}
Let $S=\Phi a$. Inside an exact equal-signal class $\mathcal A_s$, the likelihood is constant, so
\begin{equation}
 P(a\mid S=s,y)=\frac{\pi(a)}{\sum_{a'\in\mathcal A_s}\pi(a')}
 =:p_0(a\mid s).
 \label{eq:alias-conditional}
\end{equation}
These weights are independent of $y$: factorial weights under IID and uniform weights over legal count vectors
under distinct sampling. The four-user aliases in Section~\ref{sec:identifiability} are concrete instances.
An approximation that invents observation-based certainty within such a class breaks this exact constraint.

Given any positive approximation $Q$, preserve its class mass but restore the correct conditional weights:
$\bar Q(a)=Q_S(s)p_0(a\mid s)$. The KL chain rule gives
\[
 D(P\|Q)=D(P_S\|Q_S)+\E_{P_S}D(p_0(\cdot\mid S)\|Q(\cdot\mid S)),
\]
and the analogous decomposition for $D(Q\|P)$. The second term vanishes after this correction. Consequently
neither KL direction increases, and
\[
 \operatorname{TV}(P,\bar Q)=\operatorname{TV}(P_S,Q_S)\leq\operatorname{TV}(P,Q).
\]
This does not guarantee a better Bayes allocation on every observation: distributional error and the sign of a
particular decision contrast are different quantities. Finding complete alias classes may itself be expensive.
Known disjoint groups of equal-signal configurations can enforce the same local conditional rule without
claiming that all full-bank aliases have been enumerated.

For a count-tail contrast $H$, the exact signal-level observable is
$\bar H(s)=\E_{p_0(\cdot\mid s)}H(a)$, and $\E_PH=\E_{P_S}\bar H(S)$. A class with zero conditional contrast
contributes zero to that comparison's numerator. This connects identity-preserving markers with decision-focused
compression; selecting an arbitrary representative identity would not preserve the observable.

Slot permutations give a separate simplification. Averaging a slot approximation over all user permutations
preserves its count distribution and hence every multiset decision. It removes artificial slot labels but
cannot correct a wrong count law. Restoring equal-signal conditionals can change counts because different
multisets may have the same signal. These two symmetries should not be conflated.

\subsection{Exact reweighting of decision contrasts}
On common legal support, write $P=QR/\E_QR$. For any competing-copy contrast $H$,
\begin{equation}
 \E_PH-\E_QH=\frac{\operatorname{Cov}_Q(H,R)}{\E_QR}
 =\int_0^1\operatorname{Cov}_{Q_t}(H,\log R)\,dt,
 \quad Q_t\propto Q R^t.
 \label{eq:combined-sensitivity}
\end{equation}
The integral follows by differentiating a normalized finite sum. It generalizes the closing-edge identity;
computing it may still be expensive. A small derivative at $t=0$ alone is not a bound on the whole change.

If $R=R_1R_2$, corrections telescope under successively updated distributions, not under one fixed $Q$.
For independent uniform $U,V\in\{-1,1\}$, take $H=UV$, $R_1=1+U/2$, $R_2=1+V/2$. Either factor alone leaves
$\E H=0$; together they give $\E H=1/4$. Thus separate screening can miss a combined effect. Support deletion
requires an omitted-mass argument first; the positive same-support interpolation does not cover it directly.

\clearpage
\section{Exact-stage status and retained verification record}
The fixed known-$K$ model has a complete finite decoder from observation to multiset output, a proof of
minimum posterior expected PUPE and independent finite evaluations. Optimality is relative to the stated
source law and loss; noise and exact encoding ambiguity still permit errors.

Historical mathematical checks are indexed below. They are retained evidence, not a request to run another
suite or a substitute for the derivations. The 16 October separation of exact and approximate work adds no
implementation, decoder test or performance experiment.
\begin{center}\small
\begin{tabularx}{\textwidth}{@{}lY@{}}\toprule
Module suffix in \texttt{tests/} & Exact scope retained in this report\\\midrule
\texttt{polar\_posterior\_test} & Count/slot equivalence, priors, Bayes loss, conditional polar interface and pair identifiability\\
\texttt{polar\_exact\_inference\_test} & Residual/inside--outside counts, identity-preserving merging and exact decision certificates\\
\texttt{polar\_alternative\_inference\_test} & Fourier, joint elimination and complete connected expansion\\
\texttt{polar\_exact\_reductions\_test} & Sufficient states, completion bounds, exact spans and aliases\\
\texttt{polar\_exact\_computation\_test} & Compatible schedules, mixed-state obstruction and certified search\\
\texttt{polar\_exact\_reuse\_test} & Reused bounds, adverse prior masses, pair operators and rooted trees\\
\texttt{polar\_three\_user\_test} & Full K3 counts, joint parity quotients and the cycle/clique boundary\\
\texttt{polar\_synthesis\_test} & Exact recentering, alias conditionals, symmetry and reweighting identities\\\bottomrule
\end{tabularx}
\end{center}
Some preserved modules also contain checks of the removed approximation formulations. Those checks do not
extend this report's scope. Commands, finite fixtures and structural audit records live in
\path{docs/EXPERIMENT_BANK.md}. Floating checks are not outward-rounded machine certificates.

\textbf{The broad exact-derivation phase is complete.} Available constructions expose substantial exact
structure but no established tractable full-system decoder. This is not a universal impossibility theorem.
Further exact work is useful when it resolves a specific bottleneck encountered in algebraic approximation
design. A general K4 expansion or another open-ended exact survey is not a prerequisite.

\textbf{Separate approximation study.} Report 9, \path{09_polar_bayesian_approximations.tex/pdf}, owns
assumptions, their algebraic motivation, computational benefit and associated error control. The Bayes
count-tail objective remains fixed. No approximation family is selected here. Unknown $K$, fading and
changes to the source or loss require explicit extensions of the present contract.

\clearpage
\section*{Sources, provenance and implementation boundary}
This report replaces \texttt{07\_polar\_d2\_inference.tex/pdf} rather than extending its D2-first organization.
The first shared discussion motivated the common-posterior starting point:
\href{https://chatgpt.com/share/6acb8c8b-7708-83eb-a1de-175361ec21a9}{Polar decoder derivation comparison}.
The later \href{https://chatgpt.com/share/6acc145d-52dc-83eb-bc9e-bd31e185960c}{exact-inference discussion}
motivated continuing the residual recursion before choosing approximations. Its expected-count readout applies to
the old membership loss; this extension instead derives the count distributions needed by the corrected loss.
These links record discussion provenance, not authority for the mathematical claims. Equations here were derived for the
repository's encoder and sampling/metric conventions and checked as described above.

Report 8 supplies the balancing construction. The current model was checked against
\path{baselines/polar.py}, \path{baselines/odma_polar.py},
\path{framework/polar_mapping.py} and \path{benchmarks/polar_energy_geometry.py}.
The retained D2 comparison was checked against \path{framework/learned_decoders.py}.
The working reasoning and corrections are in \path{docs/POLAR_DECODING_RESEARCH_LOG.md}.
The finite certificate modules are historical implementation evidence. The 16 October revision separates
approximation design from this exact record; no encoder, decoder, training, HPC or test code is changed.

\begin{thebibliography}{9}
\setlength{\itemsep}{0.6em}
\bibitem{arikan} E. Ar{\i}kan, \emph{Channel Polarization: A Method for Constructing Capacity-Achieving Codes for
Symmetric Binary-Input Memoryless Channels}, IEEE Transactions on Information Theory, 2009.
\href{https://arxiv.org/abs/0807.3917}{arXiv:0807.3917}.
\bibitem{tal} I. Tal and A. Vardy, \emph{List Decoding of Polar Codes}, IEEE Transactions on Information Theory,
2015; preprint 2012. \href{https://arxiv.org/abs/1206.0050}{arXiv:1206.0050}.
\bibitem{wainwright} M. J. Wainwright and M. I. Jordan, \emph{Graphical Models, Exponential Families, and
Variational Inference}, 2008, especially Sec.~2.5 on exact elimination and junction trees.
\href{https://people.eecs.berkeley.edu/~jordan/sail/readings/wainwright-jordan-fnt.pdf}{Authors' manuscript}.
\bibitem{pulvirenti} E. Pulvirenti and D. Tsagkarogiannis, \emph{Cluster expansion in the canonical ensemble},
2011 preprint. \href{https://arxiv.org/abs/1105.1022}{arXiv:1105.1022}.
Cited for the classical connected-expansion context; its convergence theorem is not applied to this decoder.
\end{thebibliography}
\end{document}
